\documentclass[12pt]{article}
\usepackage[margin=0.9in]{geometry}
\usepackage[T1]{fontenc}
\usepackage{lmodern,amsmath,amssymb,amsthm,enumitem,microtype,bm}
\usepackage[colorlinks=true,linkcolor=blue,urlcolor=blue,citecolor=blue]{hyperref}
\usepackage{tikz-cd}
\renewcommand{\ker}{\operatorname{Ker}}
\newtheorem{lemma}{Lemma}
\newtheorem{theorem}{Theorem}
\newtheorem{proposition}{Proposition}
\newtheorem{corollary}{Corollary}
\newtheorem{problem}{Problem}
\newtheorem{remark}{Remark}
\newtheorem{example}{Example}
\newcommand{\Reg}{\operatorname{Reg}}
\newcommand{\RegMax}{\operatorname{RegMax}}
\newcommand{\RegSpec}{\operatorname{RegSpec}}
\newcommand{\Frac}{\mathrm{T}}
\setlist[enumerate]{itemsep=3pt,topsep=5pt}
\title{Comparable Overrings of Marot Rings II}
\author{Ba\c{s}ak Ay Saylam \& Eren Canan}
\date{}
\begin{document}
	\maketitle

\section{Introduction}

    For a commutative ring $R$, the  set of overrings of $R$ form a partially ordered set under inclusion. An overring comparable to every member of this set separates the overrings into those it contains and those which contain it. Ayache studied such overrings for integral domains in \cite{ayache}, obtaining structural results and characterizations involving valuation domains, integral closure and normal pairs. Our earlier work \cite{marot} we developed the corresponding framework for Marot rings and extended part of this theory to rings with zero divisors. Building upon our previous work, this paper focuses on comparable overrings that are integral over the base ring, and hence generalizes the results of the final section of Ayache's paper \cite{ayache} to Marot rings.

    For rings with zero divisors, regular ideals and regular localizations take the place of several constructions used for domains. A Marot ring is a ring whose regular ideals are generated by regular elements. This condition allows us to use valuation overrings and to check inclusions of regular submodules on their regular elements. A ring with a unique regular maximal ideal need not be local: additional maximal ideals may contain only zero divisors. We, therefore, formulate our results in terms of regular maximal ideals throughout.

    We first show that an integral comparable overring has at most two regular maximal ideals. When the base ring and the overring share the same unique regular maximal ideal, we characterize comparability in terms of the associated valuation overring and the intermediate residue fields. In the case of two regular maximal ideals, the integral closure is the intersection of two valuation overrings and is a Marot Prüfer ring. The intersection of the two maximal ideals determines a pullback ring whose minimal overring is the given comparable overring. This description yields to minimality criteria, restrictions on other comparable overrings, formulas for conductors and divisorial closure.

    Normal pairs provide the passage from intermediate rings in the integral closure to arbitrary overrings. We use the results of Ben Nasr and Jarboui \cite{normal-pairs} and record a simple consequence: every intermediate ring strictly below an integral comparable overring is normally closed. Finite chain conditions give another criterion for normal closedness, leading to a characterization for $FC$ Marot rings. We also obtain a criterion involving the radical of a conductor.
	
	This paper is organized as follows. \S \ref{sec-preliminaries} introduces the notation and terminology used throughout the paper. \S \ref{sec-auxiliary} presents auxiliary results concerning pullback constructions, valuation overrings and normal pairs. \S \ref{sec-integral} establishes structural results for integral comparable overrings and a comparability criterion for $FC$ Marot rings, followed by a separate treatment of Types I, II and III. It also shows that the converse of Proposition \ref{radical-conductor-comparable} does not hold in general. Finally, \S \ref{counterexamples} presents counterexamples to the conjectures posed in \cite[Problems 29 and 30]{ayache}.
\section{Preliminaries}\label{sec-preliminaries}
	
	Throughout this paper, all rings are commutative with identity, and all ring extensions preserve the identity. An element is called \emph{regular} if it is not a zero divisor. An ideal is regular if it contains a regular element. For a subset $A$ of a ring, $\Reg(A)$ denotes the set of regular elements belonging to $A$. We write $\RegSpec(R)$ and $\RegMax(R)$ for the sets of regular prime ideals and regular maximal ideals of $R$, respectively. When $R$ has a unique regular maximal ideal, we denote it by $\mathfrak{m}$.
	
	The total ring of fractions of $R$ is $T(R)=\Reg(R)^{-1}R$. An \emph{overring} of $R$ is a ring $S$ satisfying $R\subseteq S\subseteq T(R)$. For a ring extension $R\subseteq T$, a ring $S$ is called an \emph{intermediate ring} if $R\subseteq S\subseteq T$. The integral closure of $R$ in $T(R)$ is denoted by $\overline R$. For a prime ideal $P$ of $R$, the \emph{regular localization} of $R$ at $P$ is
	\[ R_{(P)}	= \left\{\frac{a}{s}:a\in R,\ s\in\Reg(R\setminus P)\right\}.\]
	
	A ring $R$ is called \emph{Marot} if every regular ideal of $R$ is generated by its regular elements. Equivalently, every regular $R$-submodule of $T(R)$ is generated by its regular elements \cite[Theorem 7.1]{huckaba}. Here an $R$-submodule of $T(R)$ is regular if it contains a regular element of $T(R)$. Consequently, inclusions between such submodules can be checked on regular elements. Every overring of a Marot ring is again Marot \cite[Corollary 7.3]{huckaba}.
	
	A Marot ring $V\neq T(V)$ is valuation ring if for every $x\in\Reg(T(V))$, either $x\in V$ or $x^{-1}\in V$ \cite[Theorem 7.7]{huckaba}. A Marot valuation ring has a unique regular maximal ideal, which we denote by $M_V$. If $R\subseteq V$, the ideal $M_V\cap R$ is called the \emph{center} of $V$ on $R$. We use the fact that the integral closure of a Marot ring in its total ring of fractions is the intersection of its Marot valuation overrings \cite[Theorem 9.3]{huckaba}. 
	
	A domain is called a \emph{discrete valuation domain} if it is a local principal ideal domain which is not a field.
	
	A ring is called \emph{Prüfer} if every finitely generated regular ideal is invertible. For a Marot ring $R$, this is equivalent to $R_{(M)}$ being a valuation ring for every $M\in\RegMax(R)$ \cite[Theorems 6.2 and 7.6]{huckaba}.
	
	Following \cite{marot}, an overring $R_o$ of $R$ is called \emph{comparable} if $R\subsetneq R_o\subsetneq T(R)$ and every overring of $R$ is comparable to $R_o$ under inclusion. A ring extension $R\subsetneq S$ is called \emph{minimal} if there is no ring strictly between $R$ and $S$. A comparable overring $R_o$ is called a \emph{minimal overring} of $R$ if $R\subsetneq R_o$ is a minimal ring extension. In this case, every overring $S\neq R$ of $R$ contains $R_o$. In particular, the minimal overring, when it exists, is unique.
	
	A regular prime ideal $P$ of $R$ is called \emph{$r$-divided} if it is comparable to every regular principal ideal of $R$ \cite[Definition 2.7]{marot}. It is called \emph{$r$-strongly prime} if $x,y\in\Reg(T(R))$ and $xy\in P$ imply $x\in P$ or $y\in P$. A ring is an \emph{$r$-pseudo-valuation ring}, or \emph{$r$-PVR}, if every regular prime ideal is $r$-strongly prime \cite{chang}. A Marot ring with unique regular maximal ideal $\mathfrak{m}$ is an $r$-PVR if and only if it has a unique Marot valuation overring with the same regular maximal ideal $\mathfrak{m}$ \cite[Theorem2.16]{marot}.\looseness=-1
	
	For $R$-submodules $A,B$ of $T(R)$, we write $(A:B)=\{x\in T(R):xB\subseteq A\}.$ For an overring $S$ of $R$ and an element $x\in T(R)$, we also use $(S:x)=\{s\in S:sx\in S\}$ for the denominator ideal of $x$ in $S$. If $S$ is an overring of $R$, then $(R:S)$ is the \emph{conductor} of $R$ in $S$. For a regular ideal $I$ of $R$, its divisorial closure is $I_v=(R:(R:I))$. The ideal $I$ is called \emph{divisorial} if $I_v=I$.
	
	A ring $R$ is called an \emph{$FC$ ring} if every chain of distinct overrings of $R$ is finite. An extension $A\subseteq B$ has the \emph{finite chain property}, or \emph{FCP}, if every chain of intermediate rings between $A$ and $B$ is finite. Thus, $R$ is an $FC$ ring precisely when $R\subseteq T(R)$ has FCP. 

	Recall that a pair $(S,T)$ is called \emph{normal} if every intermediate ring between $S$ and $T$ is integrally closed in $T$. A ring $S$ is \emph{normally closed} if, for every overring $T$ of $S$, normality of $(S,T)$ implies $T=S$. An extension $A\subseteq B$ is called a \emph{$P$-extension} if, for every $b\in B$, there exists a polynomial $f(X)=a_0+a_1X+\cdots+a_nX^n\in A[X]$ such that $f(b)=0$ and $(a_0,a_1,\ldots,a_n)=A$.
	
	If $A\subseteq B$ is integral and $(A,B)$ is normal, then $A=B$: every element of $B$ is integral over $A$ while $A$ is integrally closed in $B$. Being normally closed does not imply that the ring is integrally closed.

    Finally, let $T$ be a ring, $I$ an ideal of $T$, and let $D$ be a subring of $T/I$. If $\pi:T\longrightarrow T/I$ is the canonical quotient map, we denote the pullback ring $\pi^{-1}(D)$ by $(T;I;D)$. If $R\subseteq T$ and $D=\pi(R)$, then $\pi^{-1}(D)=R+I$ and $D\cong R/(I\cap R)$. The following commutative diagram describes the pullback ring $(T;I;D)$:
    \[
    \begin{tikzcd}[column sep=large, row sep=large]
    	(T;I;D) \arrow[r, hook] \arrow[d, two heads] & T \arrow[d, two heads, "\pi"] \\
    	D \arrow[r, hook] & T/I
    \end{tikzcd}
    \]
    The horizontal maps are inclusions, and the left vertical map is the restriction of $\pi$ to $(T;I;D)$. The diagram commutes because both paths send each $t\in(T;I;D)$ to $\pi(t)\in T/I$.

	\section{Auxiliary results}\label{sec-auxiliary}
	\subsection{Pullbacks and valuation centers}
	We begin with some auxiliary results on ideals, pullbacks and valuation overrings which are useful tools to prove our main results.
	\begin{lemma}\label{minimal-ideals}
		Let $R\subsetneq S$ be a minimal ring extension. If $I$ and $J$ are distinct ideals of $S$ and $I\cap R=J\cap R=H$, then $I\cap J=H$. Indeed, $I\cap J$ is an ideal of $R$.
	\end{lemma}
	\begin{proof}
		The proof is similar to \cite[Lemma 2.1]{gilmer-heinzer}.
	\end{proof}
	\begin{lemma}\label{comparable-ideal}
		Let $R$ be a commutative ring and $R_o$ a comparable overring of $R$. Let $T$ be an overring of $R_o$ and $J$ a proper ideal of $T$. If $J\neq J\cap R_o$, then $R_o=R+(J\cap R_o).$
	\end{lemma}
	
	\begin{proof}
		The proof is identical to that of \cite[Lemma 4]{ayache} and remains valid for commutative rings with zero divisors.
	\end{proof}
	Let $T$ be an overring of $R$, let $I$ be an ideal of $T$, and let $\pi:T\longrightarrow T/I$ be the quotient map. Since $R+I=\pi^{-1}(\pi(R))$, the ring $R+I$ is the pullback $(T;I;\pi(R))$. If $I\cap R=\mathfrak m$, then $\pi(R)\cong R/\mathfrak m$. The next lemma and its corollaries give conditions under which this pullback has a unique regular maximal ideal.
	\begin{lemma}\label{unit-pullback}
		Let $R$ be a Marot ring with unique regular maximal ideal $\mathfrak{m}$, $T$ an overring of $R$, and $I$ an ideal of $T$ such that $I\cap R=\mathfrak{m}$. Set $R_1=R+I$. Suppose that every regular element of $R_1$ outside $I$ is a unit of $T$. Then $R_1$ has unique regular maximal ideal $I$, and $R_1/I\cong R/\mathfrak{m}$.
	\end{lemma}
	\begin{proof}
		Since $I$ is an ideal of $T$, it is also an ideal of $R_1$. The second isomorphism theorem gives
		\[
		R_1/I\cong R/(I\cap R)=R/\mathfrak{m}.
		\]
		Thus $I$ is a maximal ideal of $R_1$, and it is regular since $\mathfrak{m}\subseteq I$.
		
		Suppose that $Q$ is another regular maximal ideal of $R_1$. Since $Q\nsubseteq I$, choose $q\in\Reg(Q\setminus I)$. By hypothesis, $q^{-1}\in T$. Write $q=r+a$ with $r\in R$ and $a\in I$. Then $r\notin\mathfrak{m}$; otherwise $q\in I$. Since $r+\mathfrak{m}$ is nonzero in the field $R/\mathfrak{m}$, there exists an element $v\in R$ such that $(r+\mathfrak{m})(v+\mathfrak{m})=1+\mathfrak{m}$. Therefore, we have $rv-1\in\mathfrak{m}$ and $qv-1=(rv-1)+av\in I$. As $I$ is an ideal of $T$, $q^{-1}\in T$ and $q^{-1}(qv-1)=v-q^{-1}$, it follows that $q^{-1}=v-q^{-1}(qv-1)\in R+I=R_1$. Hence, $q\in Q$, which is a contradiction. Therefore $I$ is the unique regular maximal ideal of $R_1$.
	\end{proof}

\begin{corollary}\label{two-pullback}
		Let $R$ be a Marot ring with unique regular maximal ideal $\mathfrak{m}$. Let $\{S_\lambda\}_{\lambda\in\Lambda}$ be a nonempty family of overrings of $R$, each having a unique regular maximal ideal $M_\lambda$, such that $M_\lambda\cap R=\mathfrak{m}$ for every $\lambda\in\Lambda$. Set
		\[
		T=\bigcap_{\lambda\in\Lambda}S_\lambda,\qquad
		I=\bigcap_{\lambda\in\Lambda}M_\lambda,\qquad R_1=R+I.
		\]
		Then $I$ is an ideal of $T$, and $R_1$ has unique regular maximal ideal $I$. Moreover, $M_\lambda\cap R_1=I$ for every $\lambda\in\Lambda$.
	\end{corollary}
	\begin{proof}
		Clearly $I$ is an ideal of $T$ and $I\cap R=\mathfrak{m}$. For every $\lambda\in\Lambda$, the inclusion $I\subseteq M_\lambda\cap R_1$ is clear. If $r+a\in M_\lambda\cap R_1$ with $r\in R$ and $a\in I$, then $r=(r+a)-a\in M_\lambda\cap R =\mathfrak{m}\subseteq I. $ Hence $r+a\in I$, proving $M_\lambda\cap R_1=I$ for every $\lambda\in\Lambda$. If $q=r+a\in\Reg(R_1\setminus I)$, with $r\in R$ and $a\in I$, then $r\notin\mathfrak{m}$, so $q\notin M_\lambda$ for every $\lambda$. Since $M_\lambda$ is the unique regular maximal ideal of $S_\lambda$, the regular element $q$ is a unit of each $S_\lambda$; hence $q^{-1}\in T$. The result follows from Lemma \ref{unit-pullback}.
	\end{proof}

\begin{corollary}\label{regular-pullback}
		Let $R$ be a Marot ring with unique regular maximal ideal $\mathfrak{m}$, and let $T$ be an overring of $R$. Let $I$ be a regular ideal of $T$ such that $I\cap R=\mathfrak{m}$ and $I\subseteq M$ for every $M\in\RegMax(T)$. Set $R_1=R+I$. Then $R_1$ has unique regular maximal ideal $I$, and $R_1/I\cong R/\mathfrak{m}$. Moreover, $M\cap R_1=I$ for every $M\in\RegMax(T)$.
	\end{corollary}
	\begin{proof}
		For every $M\in\RegMax(T)$, $\mathfrak{m}\subseteq M\cap R$ implies that $M\cap R=\mathfrak{m}$. If $r+a\in M\cap R_1$ with $r\in R$ and $a\in I\subseteq M$, then $r=(r+a)-a\in M\cap R=\mathfrak{m}\subseteq I$. Therefore, $r+a\in I$, and so $M\cap R_1\subseteq I$. The reverse inclusion is clear. Consequently, every regular element of $R_1$ outside $I$ lies in no regular maximal ideal of $T$ and is, therefore, a unit of $T$. The result follows from Lemma \ref{unit-pullback}.
	\end{proof}

	\begin{lemma}\label{centered}
		Let $R$ be a Marot ring and let $P\in\RegSpec(R)$. Let $\mathcal N$ be the family of all Marot valuation overrings of $R$ centered at $P$, that is,
		\[\mathcal N=\{V\mid V\text{ is a Marot valuation overring of }R\text{ and }M_V\cap R=P\}.\]
		Then the integral closure of $R_{(P)}$ in $T(R)$ is given by
		\[\overline{R_{(P)}}=\bigcap_{V\in\mathcal N}V.\]
	\end{lemma}
	\begin{proof}
		We first claim that
		\[
		\mathcal N=\{W\mid W\text{ is a Marot valuation overring of }R_{(P)}\text{ and }M_W\cap R_{(P)}=PR_{(P)}\}.
		\]
		First, let $V\in\mathcal N$. We show that $R_{(P)}\subseteq V$. Let $x=a/b\in R_{(P)}$, where $a\in R$ and $b\in\Reg(R\setminus P)$. Since $b\in\Reg(R\setminus P)$ and $M_V\cap R=P$, $b$ is a unit of $V$. Since $a\in R\subseteq V$, it follows that $x=ab^{-1}\in V$. Thus $R_{(P)}\subseteq V$. Moreover, one can easily show that $M_V\cap R_{(P)}=PR_{(P)}$. Indeed, for $a/b\in R_{(P)}$ with $a\in R$ and $b\in\Reg(R\setminus P)$, we have
		\[
		\frac{a}{b}\in M_V\Longleftrightarrow a\in M_V\Longleftrightarrow a\in P\Longleftrightarrow\frac{a}{b}\in PR_{(P)}.
		\]
		Thus, $V$ is a Marot valuation overring of $R_{(P)}$ centered at $PR_{(P)}$. Conversely, let $W$ be a Marot valuation overring of $R_{(P)}$ such that $M_W\cap R_{(P)}=PR_{(P)}$. Since $R\subseteq R_{(P)}\subseteq W$, $W$ is also a Marot valuation overring of $R$. Furthermore, we have $M_W\cap R=(M_W\cap R_{(P)})\cap R=PR_{(P)}\cap R=P$. Therefore, $W\in\mathcal N$, establishing the claimed equality of the two families of valuation overrings.
		
		Next, we show that every Marot valuation overring $S\neq T(R)$ of $R_{(P)}$ contains some $W\in\mathcal N$. Let $Q=M_S\cap R_{(P)}$. By \cite[Theorem 2.4]{marot}, $Q$ is a regular prime ideal of $R_{(P)}$, so $Q\subseteq PR_{(P)}$. Since $Q$ is regular and $PR_{(P)}$ is the unique regular maximal ideal of $R_{(P)}$, the quotient $R_{(P)}/Q$ is a local domain with maximal ideal $PR_{(P)}/Q$. Let $k=S/M_S$, and let $\pi:S\to k$ be the canonical surjection; identify $R_{(P)}/Q$ with its image under $\pi$, a subring of $k$. By \cite[Theorem 5.2]{huckaba}, there exists a valuation domain $U$ of $k$ such that
		\[
		R_{(P)}/Q\subseteq U\subseteq k,\qquad \mathfrak{m}_U\cap(R_{(P)}/Q)=PR_{(P)}/Q,
		\]
		where $\mathfrak{m}_U$ is the maximal ideal of $U$; that is, $U$ dominates $R_{(P)}/Q$.
		
		Set $W=\pi^{-1}(U)=\{s\in S\mid s+M_S\in U\}$. Our aim is to show that $W$ is a Marot valuation overring of $R_{(P)}$ with $M_W\cap R_{(P)}=PR_{(P)}$, i.e., $W\in\mathcal N$, and $W\subseteq S$. Let $x\in\Reg(T(R))$. Since $S$ is a Marot valuation ring, $x\in S$ or $x^{-1}\in S$; say $x\in S$ (the other case is symmetric). Every element of $M_S$ has image $0\in U$ under $\pi$, so the definition $W=\pi^{-1}(U)$ gives $M_S\subseteq W$. If $\pi(x)=0$, then $x\in M_S\subseteq W$. If $\pi(x)\neq0$, then $x\notin M_S$. Since $x$ is regular and $M_S$ is the unique regular maximal ideal of $S$, $x$ is a unit of $S$. Hence, $\pi(x^{-1})=\pi(x)^{-1}$. As $U$ is a valuation domain with quotient field $k$, either $\pi(x)\in U$ or $\pi(x)^{-1}\in U$. Consequently, either $x\in W$ or $x^{-1}\in W$. On the other hand, the unique regular maximal ideal of $W$ is $M_W=\pi^{-1}(\mathfrak{m}_U)$, and
		\[
		M_W\cap R_{(P)}=\pi^{-1}(\mathfrak{m}_U)\cap R_{(P)}= \pi|_{R_{(P)}}^{-1} \bigl( \mathfrak{m}_U \cap(R_{(P)}/Q)\bigr)=\pi|_{R_{(P)}}^{-1}(PR_{(P)}/Q)=PR_{(P)}.
		\]
		Finally, since every Marot valuation overring $S$ of $R_{(P)}$ contains a Marot valuation overring $W$ centered at $PR_{(P)}$ (i.e., $W\in\mathcal N$), the intersection over all valuation overrings of $R_{(P)}$ reduces precisely to the intersection over the family $\mathcal N$. Since the integral closure $\overline{R_{(P)}}$ of $R_{(P)}$ in $T(R)$ is the intersection of all Marot valuation overrings of $R_{(P)}$ \cite[Theorem 9.3]{huckaba}, we obtain
		\[
		\overline{R_{(P)}}=\bigcap_{V\in\mathcal N}V,
		\]
		completing the proof.
	\end{proof}

	\subsection{Normal pairs and finite chains}
	Recall that a pair of rings $(S,T)$ with $S\subseteq T$ is called a \emph{normal pair} if every intermediate ring $H$, with $S\subseteq H\subseteq T$, is integrally closed in $T$.
	\begin{lemma}\label{normal-intersection}
		Let $R$ be a ring such that $\overline R$ is Prüfer. For every overring $T$ of $R$, the pair $(T\cap\overline R,T)$ is normal.
	\end{lemma}
	
	\begin{proof}
		Put $S=T\cap\overline R$. Since $\overline R$ is Prüfer, the pair $(\overline R,\Frac(R))$ is normal. By \cite[Theorem 1]{normal-pairs}, $\overline R\subseteq \Frac(R)$ is a $P$-extension. On the other hand, for every $x\in\overline R$, there exist $r_0,\ldots,r_{n-1}\in R$ such that $x^n+r_{n-1}x^{n-1}+\cdots+r_0=0$. Since $(r_0,\ldots,r_{n-1},1)=R$, the integral extension $R\subseteq\overline R$ is also a $P$-extension. Thus, \cite[Theorem 2]{normal-pairs} shows that $R\subseteq \Frac(R)$ is a $P$-extension. For each $t\in T$, there exists a polynomial $f(X)=r_0+r_1X+\cdots+r_nX^n\in R[X]$ such that $f(t)=0$ and $(r_0,\ldots,r_n)=R$. Since $R\subseteq S$, $f(X)\in S[X]$, and $(r_0,\ldots,r_n)S=S$. Hence $S\subseteq T$ is a $P$-extension. We next show that $S$ is integrally closed in $T$. Let $t\in T$ be integral over $S$. Since $R\subseteq S\subseteq\overline R$, the extension $R\subseteq S$ is integral. By transitivity of integrality, $t$ is integral over $R$, so $t\in\overline R$. Thus $t\in T\cap\overline R=S$. Hence, $S$ is integrally closed in $T$. Therefore, $S\subseteq T$ is a $P$-extension and $S$ is integrally closed in $T$. By \cite[Theorem 1]{normal-pairs}, $(S,T)$ is a normal pair.
	\end{proof}

The following observation isolates the normal-closedness argument used in the proof of \cite[Theorem 24]{ayache}. It requires neither the Marot condition nor a hypothesis on the number of regular maximal ideals.

\begin{lemma}\label{integral-comparable-normally-closed}
    Let $R$ be a ring, and let $R_o$ be a comparable overring of $R$ which is integral over $R$. Then every intermediate ring $S$ with $R\subseteq S\subsetneq R_o$ is normally closed. In particular, $R$ is normally closed.
\end{lemma}
\begin{proof}
    Let $T$ be an overring of $S$ such that $(S,T)$ is normal. On the other hand, $T$ and $R_o$ are comparable. Suppose that $R_o\subseteq T$, and let $x\in R_o$. Since $R_o$ is integral over $R$, there exists a monic polynomial $f(X)\in R[X]\subseteq S[X]$ such that $f(x)=0$. Hence, $x$ is integral over $S$. Since $x\in R_o\subseteq T$ and $S$ is integrally closed in $T$, we obtain $x\in S$. Thus, $R_o\subseteq S$, contradicting to $S\subsetneq R_o$. Therefore, $T\subsetneq R_o$. Now, $T$ is integral over $S$, so, the same integral closedness gives $T=S$. Hence, $S$ is normally closed.
\end{proof}
	We now turn to finite chain conditions. The $P$-extension assertion in the next lemma is also a special case of \cite[Proposition 2.4]{dobbs-jarboui}.
	\begin{lemma}\label{fc-prufer-closure}
			Let $R$ be an $FC$ ring. Then, $R\subseteq T(R)$ is a $P$-extension, and $\overline R$ is Prüfer.
	\end{lemma}
	\begin{proof}
		Since $R$ is an $FC$ ring, both of the extensions $R\subseteq T(R)$ and $\overline R\subseteq T(R)$ satisfy FCP. By \cite[Proposition 2.4]{dobbs-jarboui}, both extensions are $P$-extensions. Moreover, $\overline R$ is integrally closed in $T(R)$. Thus \cite[Theorem 1]{normal-pairs} implies that $(\overline R,T(R))$ is a normal pair. It follows that $\overline R$ is Prüfer.
	\end{proof}
	The following lemma shows that, for regular elements in overrings, the quasilocal hypothesis in \cite[Lemma 1]{normal-pairs} can be replaced by the existence of a unique regular maximal ideal.
	\begin{lemma}\label{normal-regular-inverse}
		Let $R$ be a ring with unique regular maximal ideal $\mathfrak{m}$, and let $T$ be an overring of $R$. If $(R,T)$ is normal, then $x^{-1}\in\mathfrak{m}$ for every $x\in\Reg(T\setminus R)$.
	\end{lemma}
	\begin{proof}
		First, observe that for every $z\in T(R)\setminus R$, the ideal $(R:z)=\{r\in R:rz\in R\}$ is proper and regular. Indeed, it is proper because $1\notin(R:z)$. Writing $z=a/b$, where $a\in R$ and $b\in\Reg(R)$, gives $bz=a\in R$, so $b\in(R:z)$. Since $\mathfrak{m}$ is the unique regular maximal ideal of $R$, we obtain $(R:z)\subseteq\mathfrak{m}.$
		
		Let $x\in\Reg(T\setminus R)$, and suppose that $x^{-1}\notin R$. Since $(R,T)$ is normal, $R$ is integrally closed in $T$, and \cite[Theorem 1]{normal-pairs} shows that $R\subseteq T$ is a $P$-extension. Consequently, there exists a polynomial with root  $x$ whose coefficients are not all in $\mathfrak{m}$. Among all such polynomials, choose one of smallest degree, say $f(X)=a_0+a_1X+\cdots+a_nX^n$ where $ n\geq 1.$ Set $y=a_nx\in T$. Multiplying $f(x)=0$ by
		$a_n^{\,n-1}$ gives
		\[
		y^n+\sum_{j=0}^{n-1}
		a_j a_n^{\,n-1-j}y^j=0.
		\]
		Thus, $y$ is integral over $R$, and hence $y\in R$. Since $a_nx=y\in R$, we have $a_n\in(R:x)\subseteq\mathfrak{m}$. Moreover, $yx^{-1}=a_n\in R$, so $y\in(R:x^{-1})\subseteq\mathfrak{m}$. Now define $$g(X)=f(X)-a_nX^n+yX^{n-1}\in R[X]. $$ Then $$g(x)=f(x)-a_nx^n+yx^{n-1}=a_0+a_1x+\cdots+(a_{n-1}+y)x^{n-1}=0. $$ Since the coefficients of $f$ are not all in $\mathfrak{m}$ and $a_n\in\mathfrak{m}$, at least one of $a_0,\ldots,a_{n-1}$ lies outside $\mathfrak{m}$. In $g(X)$, the coefficients of the terms of degree less than $n-1$ remain unchanged, while $a_{n-1}$ is replaced by $a_{n-1}+y$. Since $y\in\mathfrak{m}$, this replacement cannot move a coefficient outside $\mathfrak{m}$ into $\mathfrak{m}$. Thus the coefficients of $g(X)$ are not all in $\mathfrak{m}$. But $\deg g(X)\leq n-1$, contradicting to the choice of $f(X)$. Therefore, $x^{-1}\in R$. Finally, $x^{-1}x=1\in R$ gives $x^{-1}\in(R:x)\subseteq\mathfrak{m}$.
	\end{proof}
	
	\begin{corollary}\label{rpvr-normally-closed}
		Let $R$ be a Marot $r$-PVR with unique regular maximal ideal $\mathfrak m$. Then $R$ is normally closed if and only if $R$ is not a valuation ring.
	\end{corollary}
	
	\begin{proof}
		Suppose first that $R$ is not a valuation ring. By \cite[Theorem 2.16]{marot}, $R$ has an associated Marot valuation overring $V$ with regular maximal ideal $\mathfrak m$. Since $R$ is not a valuation ring, $R\subsetneq V$. Let $T$ be an overring of $R$ such that $(R,T)$ is normal. Suppose, for a contradiction, that $R\subsetneq T$. By \cite[Theorem 3.1]{marot}, $T$ and $V$ are comparable under inclusion. Since $T$ and $V$ are comparable, their intersection equals to one of them. As both properly contain $R$, we have $R\subsetneq T\cap V$. Let $x\in\Reg\bigl((T\cap V)\setminus R\bigr).$ Lemma \ref{normal-regular-inverse} implies that $x^{-1}\in\mathfrak m$. Since $x\in V$ and $\mathfrak m$ is an ideal of $V$, we obtain $1=xx^{-1}\in\mathfrak m,$ a contradiction. Therefore, $T=R$, and hence, $R$ is normally closed.\looseness=-1
		
		Conversely, suppose that $R$ is a valuation ring. Then, $R$ is Prüfer, so $(R,T(R))$ is a normal pair. Choose $a\in\Reg(\mathfrak m)$. Since $a$ is not a unit of $R$, we have $a^{-1}\in T(R)\setminus R$. Thus, $R\subsetneq T(R)$ implying that $R$ is not normally closed.
	\end{proof}
	
	Lemma \ref{integral-comparable-normally-closed} shows that a ring with an integral comparable overring is normally closed. The next lemma gives another sufficient condition: a Marot ring $R$ with unique regular maximal ideal is normally closed whenever it has a proper integral overring $T$ such that $R\subsetneq T$ satisfies FCP. Here, $T$ need not be comparable.

	\begin{lemma}\label{fc-integral-normally-closed}
		Let $R$ be a Marot ring with unique regular maximal ideal $\mathfrak{m}$. Suppose that $R$ has a proper integral overring $T$ such that every chain of intermediate rings between $R$ and $T$ is finite. Then, $R$ is normally closed.
	\end{lemma}
	
	\begin{proof}		
		Let $H$ be an overring of $R$ such that $(R,H)$ is normal. Suppose, for a contradiction, that $H\neq R$. So, there exists an element $h\in\Reg(H\setminus R)$. By Lemma \ref{normal-regular-inverse}, $h^{-1}\in\mathfrak{m}.$ Set $a:=h^{-1}$. The ascending chain condition on intermediate rings gives $T=R[t_1,\ldots,t_q]$. Since each $t_i$ is integral over $R$, \cite[Corollary 5.2]{atiyah} shows that $T$ is a finitely generated $R$-module. Choose a common regular denominator $d\in\Reg(R)$ such that $dT\subseteq R$. On the other hand, for every $k\geq 1$, $R+a^kT$ is an intermediate ring between $R$ and $T$. Now, consider the descending chain $R+aT\supseteq R+a^2T \supseteq R+a^3T\supseteq\cdots.$ Since every chain of intermediate rings between $R$ and $T$ is finite, there exists $n\geq 1$ such that
		\[
		R+a^nT=R+a^{n+1}T.
		\]
		Now, set $F=(R+a^nT)/R,$ where the quotient is taken as an $R$-module. $T$ is a finitely generated $R$-module, so is $a^nT$. Furthermore, $R+a^nT$ is finitely generated, and so is $F$.
		
		We next show that $F=aF$. Since $a\in R$, we have $aF\subseteq F$. Conversely, let $a^nt+R\in F$. The equality of the two intermediate rings above gives $a^nt=r+a^{n+1}u $ for some $r\in R$ and $u\in T$. Hence, $a^nt+R=a^{n+1}u+R =a\bigl(a^nu+R\bigr)\in aF. $ Thus, $F=aF$.
		
		We next show that $dF=0$. Since $dT\subseteq R$ and $a\in R$, we have $a^n(dt)\in R$ for every $t\in T$. Consequently, $d\bigl(a^nt+R\bigr) =a^n(dt)+R=0+R.$ Therefore, $dF=0$.
		
		We next show that $F=0.$ Since $F$ is finitely generated and $(aR)F=aF=F$, \cite[Corollary 2.5]{atiyah} applied to the ideal $aR$, there exists an element $c\in R$ such that $c-1\in aR$ and $cF=0$. Since $c-1\in aR\subseteq\mathfrak{m}$, we have $c\notin\mathfrak{m}$; otherwise $1=c-(c-1)\in\mathfrak{m}$, a contradiction. The ideal $(d,c)$ is regular and is not contained in the unique regular maximal ideal $\mathfrak{m}$; hence $(d,c)=R$. Therefore, $$F=(d,c)F=dF+cF=0.$$ Consequently, $R+a^nT=R$, so $a^nT\subseteq R$. Thus, for every $t\in T$, both $h^n$ and $a^nt$ belong to $H$ since $h\in H$ and $a^nT\subseteq R\subseteq H$. Therefore, $t=h^n(a^nt)\in H. $ Thus, $T\subseteq H$. Since $T$ is integral over $R$ and $R$ is integrally closed in $H$, we obtain $T=R$, contradicting to $R\subsetneq T$. Consequently, $H=R$, so $R$ is normally closed.
	\end{proof}
	
	\begin{corollary}\label{fc-normally-closed-characterization}
		Let $R$ be an $FC$ Marot ring with unique regular maximal ideal $\mathfrak m$. The following conditions are equivalent:
		\begin{enumerate}[label=(\roman*)]
			\item $R$ is normally closed.
			\item $R\neq\overline R$.
			\item $R$ is not Prüfer.
		\end{enumerate}
	\end{corollary}

	\begin{proof}
		\emph{(ii)$\Rightarrow$(i).} Since $R$ is an $FC$ ring, the proper integral extension $R\subsetneq\overline R$ satisfies FCP. Lemma \ref{fc-integral-normally-closed} imlies that $R$ is normally closed. 
	
		\emph{(i)$\Rightarrow$(ii).} Suppose that $R=\overline R$. By Lemma \ref{fc-prufer-closure}, $R$ is Prüfer. Hence $(R,T(R))$ is a normal pair. Choose $a\in\Reg(\mathfrak m)$. Since $a$ is not a unit of $R$, we have $a^{-1}\in T(R)\setminus R$. Thus, $R\subsetneq T(R)$, contradicting to the normal closedness of $R$.
	
		\emph{(ii)$\Leftrightarrow$(iii).} A Prüfer ring is integrally closed. Conversely, if $R$ is integrally closed, then $R=\overline R$ is Prüfer by Lemma \ref{fc-prufer-closure}.
	\end{proof}

	\section{The Case Where $R_o\subseteq\overline R$}\label{sec-integral}
	We prove that an integral comparable overring of a Marot ring cannot have more than two regular maximal ideals.
	\begin{theorem}\label{at-most-two}
		If $R_o$ is a comparable overring of a Marot ring $R$ and $R_o\subseteq\overline R$, then $R_o$ has at most two regular maximal ideals.
	\end{theorem}
	\begin{proof}
		By \cite[Proposition 3.3]{marot}, $R$ has a unique regular maximal ideal $\mathfrak{m}$. Suppose $P_o,Q_o$ are distinct regular maximal ideals of $R_o$. Since $R_o$ is integral over $R$, $P_o\cap R=Q_o\cap R=\mathfrak{m}$. Choose valuation overrings $V_1,V_2$ of $R_o$ centered at $P_o,Q_o$, respectively, with regular maximal ideals $M_1,M_2$. Set
		\[
		T=V_1\cap V_2,\qquad I=M_1\cap M_2,\qquad R_1=R+I.
		\]
		Corollary \ref{two-pullback} implies that $R_1$ has unique regular maximal ideal $I$ and that $M_i\cap R_1=I$ for $i=1,2$.
		
		Now, suppose that $R_o\subseteq R_1$. Since $M_i\cap R_1=I$, we obtain $P_o=M_1\cap R_o=(M_1\cap R_1)\cap R_o=I\cap R_o$, and similarly, $Q_o=M_2\cap R_o=(M_2\cap R_1)\cap R_o=I\cap R_o$. Hence, $P_o=Q_o$, contradicting to our choice of two distinct ideals. Comparability, therefore, gives $R_1\subsetneq R_o$. Since $I\subseteq R_1\subseteq R_o$, we have $P_o\cap Q_o=(M_1\cap R_o)\cap(M_2\cap R_o)=(M_1\cap M_2)\cap R_o=I$. The extension $R_1\subseteq R_o$ is integral. So, every regular maximal ideal $N$ of $R_o$ contracts to $I$ in $R_1$ and contains $P_o\cap Q_o$. Since $P_oQ_o=P_o\cap Q_o\subseteq N$ and $N$ is prime, by \cite[Proposition 1.11]{atiyah}, we have either $P_o\subseteq N$ or $Q_o\subseteq N$. Since $P_o$ and $Q_o$ are maximal ideals, either $N=P_o$ or $N=Q_o$.
	\end{proof}

\begin{remark}\label{minimal-without-marot}
	
		In the minimal case, neither the Marot hypothesis nor comparability is needed. Let $R$ be a ring with unique regular maximal ideal $\mathfrak{m}$, and let $R_o$ be an integral overring such that $R\subsetneq R_o$ is a minimal ring extension. Then $R_o$ has at most two regular maximal ideals. Indeed, suppose that $P_o,Q_o,N$ are three distinct regular maximal ideals of $R_o$. Since $R_o$ is an integral overring of $R$, we have $P_o\cap R=Q_o\cap R=N\cap R=\mathfrak{m}.$ By Lemma \ref{minimal-ideals}, $P_o\cap Q_o=\mathfrak{m}$. Hence, $P_oQ_o=P_o\cap Q_o=\mathfrak{m}\subseteq N. $ Since $N$ is prime, $P_o\subseteq N$ or $Q_o\subseteq N$. Hence, $P_o=N$ or $Q_o=N$, a contradiction. 
	\end{remark}

As a result, there are three types of comparable overrings $R_o$ of a Marot ring $R$ with unique regular maximal ideal $\mathfrak m$ such that $R_o\subseteq\overline R$:
\begin{enumerate}[label=\arabic*.]
    \item Type I: $R_o$ has unique regular maximal ideal $Q_o=\mathfrak m$.
    \item Type II: $R_o$ has two regular maximal ideals $P_o$ and $Q_o$.
    \item Type III: $R_o$ has unique regular maximal ideal $Q_o$ such that $\mathfrak m\subsetneq Q_o$.
\end{enumerate}

Before treating these types separately, we record a minimality criterion valid whenever the integral closure is Prüfer. It extends \cite[Corollary 28]{ayache} and will allow us to avoid repeating the passage from integral intermediate rings to arbitrary overrings.

\begin{corollary}\label{integral-minimal-characterization}
		Let $R$ be a Marot ring with unique regular maximal ideal $\mathfrak{m}$ such that $\overline R$ is Prüfer. Let $R_o$ be a proper overring of $R$ contained in $\overline R$. The following conditions are equivalent:
		\begin{enumerate}[label=(\roman*)]
			\item $R_o$ is the minimal overring of $R$.
			\item Every intermediate ring $S$ with $R\subsetneq S\subseteq\overline R$ contains $R_o$.
		\end{enumerate}
	\end{corollary}
	
	\begin{proof}
		\emph{(i)$\Rightarrow$(ii).} By comparability and minimality, every proper overring of $R$ contains $R_o$.
		
		\emph{(ii)$\Rightarrow$(i).} If $R\subsetneq S\subseteq R_o$, then (ii) gives $R_o\subseteq S$, so $S=R_o$. Thus, $R\subsetneq R_o$ is a minimal ring extension. In particular, every chain of intermediate rings between $R$ and $R_o$ is finite. Since $R_o$ is integral over $R$, Lemma \ref{fc-integral-normally-closed} implies that $R$ is normally closed.
		
		Let $T$ be any proper overring of $R$, and set $S=T\cap\overline R$. By Lemma \ref{normal-intersection}, $(S,T)$ is normal. If $S=R$, normal closedness would give $T=R$, contradicting to the choice of $T$. Therefore, $R\subsetneq S\subseteq\overline R.$ By (ii), we obtain $R_o\subseteq S\subseteq T$. Hence, every proper overring of $R$ contains $R_o$. Together with the minimality proved above, this shows that $R_o$ is the minimal overring of $R$.
	\end{proof}

The following criterion applies to all three types. For $FC$ Marot rings, Lemma \ref{fc-integral-normally-closed} allows us to replace the normal-closedness hypothesis by uniqueness of the regular maximal ideal. The following theorem extends \cite[Theorem 26]{ayache}.

\begin{theorem}\label{fc-comparable-characterization}
	Let $R$ be an $FC$ Marot ring with unique regular
	maximal ideal $\mathfrak{m}$, and let $R_o$ be a
	proper overring of $R$ such that
	$R_o\subseteq\overline R$.
	Then $R_o$ is a comparable overring of $R$
	if and only if:
	\begin{enumerate}[label=(\roman*)]
		\item every intermediate ring $S$ between
		$R$ and $\overline R$ is comparable to $R_o$
		under inclusion;
		\item every intermediate ring $S$ with
		$R\subseteq S\subsetneq R_o$ has a unique
		regular maximal ideal.
	\end{enumerate}
\end{theorem}
\begin{proof}
	Suppose that $R_o$ is comparable. Assertion (i) follows immediately. Let $R\subseteq S\subsetneq R_o$. Every overring of $S$ is an overring of $R$, so $R_o$ is a comparable overring of $S$. Moreover, $R_o$ is integral over $S$. By \cite[Proposition 3.3]{marot}, $S$ has a unique regular maximal ideal. This proves (ii).
	
	Conversely, suppose that (i) and (ii) hold. By Lemma \ref{fc-prufer-closure}, $\overline R$ is Prüfer. Let $T$ be any overring of $R$, and put $S=T\cap\overline R$. By (i), either $R_o\subseteq S$, in which case $R_o\subseteq T$, or $S\subsetneq R_o$.
	
	Suppose that $S\subsetneq R_o$. By (ii), $S$ has a unique regular maximal ideal. The extension $S\subsetneq R_o$ is integral, and every chain of intermediate rings between $S$ and $R_o$ is finite because $R$ is an $FC$ ring. Since $S$ is Marot, Lemma \ref{fc-integral-normally-closed} shows that $S$ is normally closed. On the other hand, Lemma \ref{normal-intersection} implies that $(S,T)$ is normal. Therefore, $T=S\subsetneq R_o$. Thus, every overring of $R$ is comparable to $R_o$.\looseness=-1
\end{proof}

\subsection{Type I}\label{subsec-type-i}
	We first consider the case in which $R$ and $R_o$ share the same unique regular maximal ideal $\mathfrak m$. Ayache uses a criterion for valuation domains that is stated for integral domains. To follow the same argument here, we would need an additional criterion for valuation rings with zero divisors. The following theorem gives the Marot counterpart of \cite[Theorem 20]{ayache}; in particular, the associated valuation overring is the integral closure of $R$. The following theorem generalizes \cite[Theorem 20]{ayache} to Marot rings. Our proof uses residue fields and monic polynomials to identify the associated valuation overring with $\overline R$. Comparability shows that its residue field is algebraic over $R/\mathfrak m$. We then lift monic polynomial relations to show that every element of the associated valuation overring is integral over $R$. This gives a direct proof that it equals $\overline R$.

	\begin{theorem}\label{comparable-same-maximal}
		Let $R$ be a Marot ring with unique regular maximal ideal $\mathfrak{m}$, and let $R_o$ be a proper overring of $R$ such that $R_o\subseteq\overline R$ and $R_o$ has $\mathfrak{m}$ as its unique regular maximal ideal. Then $R_o$ is a comparable overring of $R$ if and only if:
		\begin{enumerate}[label=(\roman*)]
			\item $R$ is an $r$-PVR with associated Marot valuation overring $\overline R$;
			\item every intermediate field between $R/\mathfrak{m}$ and $\overline R/\mathfrak{m}$ is comparable to $R_o/\mathfrak{m}$ under inclusion.
		\end{enumerate}
	\end{theorem}
	\begin{proof}
		Suppose that $R_o$ is a comparable overring of $R$. By \cite[Corollary 3.5]{marot}, $R$ is an $r$-PVR. By \cite[Theorem 2.16]{marot}, $R$ has a unique associated Marot valuation overring $V$ with unique regular maximal ideal $\mathfrak{m}$. Since $V$ is integrally closed, $R_o\subseteq\overline R\subseteq V$. We show that $V=\overline R$.
		
		Put $k=R/\mathfrak{m}$, $K=R_o/\mathfrak{m}$ and $F=V/\mathfrak{m}$. These are fields with $k\subseteq K\subseteq F$. Moreover, $k\subsetneq K$. Otherwise, every class $a+\mathfrak{m}\in K$ would equal to $r+\mathfrak{m}$ for some $r\in R$ since $K=k=R/\mathfrak{m}$. Thus, $a-r\in\mathfrak{m}\subseteq R$, and hence $a=r+(a-r)\in R$, forcing $R_o=R$.
		
		Let $\pi:V\to F$ be the canonical projection. Suppose that some $z\in F$ is transcendental over $k$, and set $T=\pi^{-1}(k[z])$. Then $T$ is an overring of $R$, and $T/\mathfrak{m}=k[z]$. Comparability of $R_o$ and $T$ gives $k[z]\subseteq K$ or $ K\subseteq k[z].$ The first inclusion is impossible because $K/k$ is algebraic while $z$ is transcendental. Indeed, reducing a monic equation for any $a\in R_o$ modulo $\mathfrak m$ gives a monic polynomial over $k$ having $a+\mathfrak m\in K$ as a root. For the second, every nonzero element of the field $K$ would be a unit of $k[z]$. Since $z$ is transcendental, the units of $k[z]$ are precisely $k^\times$, the set of nonzero element of $k$. Indeed, if $p(z)q(z)=1$ for $p,q\in k[X]$, then $z$ being transcendental over $k$ implies $p(X)q(X)=1$, so, both $p$ and $q$ are nonzero constant polynomials. Thus, every nonzero element of $K$ belongs to $k$, so $K\subseteq k$; together with $k\subsetneq K$, this gives a contradiction. Hence, $F/k$ is algebraic.
		
		Now, let $v\in V$. Since $F/k$ is algebraic, choose a monic polynomial $g\in k[X]$ such that $g(v+\mathfrak{m})=0$ and lift its coefficients to a monic polynomial $G\in R[X]$. Then, $c:=G(v)\in\mathfrak{m}\subseteq R$. Define $h(X)=G(X)-c\in R[X]$. Then, $v$ is a root of the monic polynomial $h(X)$. Thus every element of $V$ is integral over $R$, so $V\subseteq\overline R$. Together with the reverse inclusion, this gives $V=\overline R$.
		
		Let $D$ be any intermediate field between $R/\mathfrak{m}$ and $\overline R/\mathfrak{m}$. Under the canonical projection $\pi:\overline R\to\overline R/\mathfrak{m}$, the preimage $T=\pi^{-1}(D)$ is an intermediate ring between $R$ and $\overline R$, with $T/\mathfrak{m}=D$. Comparability of $R_o$ gives $T\subseteq R_o$ or $R_o\subseteq T$. Passing to modulo $\mathfrak{m}$ proves (ii).
		
		Conversely, suppose that (i) and (ii) hold, and let $T$ be any overring of $R$. By \cite[Theorem 3.1]{marot}, $T$ and $\overline R$ are comparable under inclusion. If $\overline R\subseteq T$, then $R_o\subseteq T$. Otherwise, $T\subsetneq\overline R$. Since $\mathfrak{m}\subseteq R\subseteq T\subseteq\overline R$ and $\mathfrak{m}$ is a prime ideal of $\overline R$, it is also a prime ideal of $T$. Thus, $T/\mathfrak{m}$ is an intermediate domain between $R/\mathfrak{m}$ and $\overline R/\mathfrak{m}$. Since $\overline R/\mathfrak m$ is algebraic over the field $R/\mathfrak m$, so is its subdomain $T/\mathfrak m$. Hence, $T/\mathfrak m$ is a field. By (ii), it is comparable to $R_o/\mathfrak{m}$. Since both rings contain $\mathfrak{m}$, lifting the inclusion gives $T\subseteq R_o$ or $R_o\subseteq T$. Thus, $R_o$ is a comparable overring of $R$.\looseness=-1
	\end{proof}

The following corollary characterizes when $R_o$ is the minimal overring of $R$.

	\begin{corollary}\label{minimal-same-maximal}
		Let $R$ be a Marot ring with unique regular maximal ideal $\mathfrak{m}$, and let $R_o$ be a proper overring of $R$ such that $R_o\subseteq\overline R$ and $R_o$ has $\mathfrak{m}$ as its unique regular maximal ideal. Then $R_o$ is the minimal overring of $R$ if and only if:
		\begin{enumerate}[label=(\roman*)]
			\item $R$ is an $r$-PVR with associated Marot valuation overring $\overline R$;
			\item every intermediate field $D$ between $R/\mathfrak{m}$ and $\overline R/\mathfrak{m}$ different from $R/\mathfrak{m}$ contains $R_o/\mathfrak{m}$.
		\end{enumerate}
	\end{corollary}

\subsection{Type II}\label{subsec-type-ii}
	We now consider comparable overrings with two regular maximal ideals. The next proposition gives a structural result needed to generalize \cite[Proposition 23]{ayache}. Ayache uses \cite[Proposition 2.5]{gilmer-heinzer}. Gilmer and Heinzer use the fact that each valuation domain
	is a localization of their finite intersection. Here we prove directly that the regular localizations of
	$\overline R$ at $N_1$ and $N_2$ are $V_1$ and $V_2$, respectively. Here we first prove that $M_1\cap M_2=\mathfrak m$. We then use monic polynomials to show that every element of $V_1\cap V_2$ is integral over $R$, which gives $V_1\cap V_2=\overline R$. Next, we prove that there is only one Marot valuation overring at each of the two centers. Finally, we prove that $(\overline R)_{(N_i)}=V_i$ for $i=1,2$.
	\begin{proposition}\label{minimal-two}
		Let $R$ be a Marot ring with unique regular maximal ideal $\mathfrak{m}$, and let $R_o$ be its minimal overring. Suppose $R_o$ has exactly two regular maximal ideals $P_o,Q_o$. Let $V_1$ and $V_2$ be valuation overrings of $R_o$ centered at $P_o$ and $Q_o$, respectively, with their unique regular maximal ideals $M_1$ and $M_2$. Then $\overline R=V_1\cap V_2$ and $M_1\cap M_2=\mathfrak{m}$. Moreover, the only Marot valuation overrings of $R$ centered at $\mathfrak{m}$ are $V_1$ and $V_2$. The ideals $N_i=M_i\cap\overline R$ are precisely the two regular maximal ideals of $\overline R$, and $\overline R_{(N_i)}=V_i$ for $i=1,2$. In particular, $\overline R$ is a Marot Prüfer ring.
	\end{proposition}
	\begin{proof}
		
		Since $R_o$ is a minimal overring of $R$, it is comparable.	As $R_o$ has exactly two regular maximal ideals, it is not a Marot valuation ring. Therefore, \cite[Proposition 3.2]{marot} implies that $R_o$ is integral over $R$. Thus $P_o\cap R=Q_o\cap R=\mathfrak{m}$. Lemma \ref{minimal-ideals} gives $P_o\cap Q_o=\mathfrak{m}$. 
		
		\smallskip We first claim that $M_1\cap M_2=\mathfrak{m}$. Clearly, $\mathfrak{m}\subseteq M_1\cap M_2$. If $x\in(M_1\cap M_2)\setminus \mathfrak{m}$, then $x\notin R$ and $R_o\subseteq R[x]$. For any $t\in R_o$, write $t=r_0+r_1x+\cdots+r_nx^n$ with $r_i\in R$. Then $t-r_0\in(M_1\cap M_2)\cap R_o=P_o\cap Q_o=\mathfrak{m}\subseteq R$. This would imply $R_o\subseteq R$, a contradiction. Hence $M_1\cap M_2=\mathfrak{m}$.
		
		\smallskip Now, let us show that $\overline R=V_1\cap V_2$. Set $T=V_1\cap V_2$ and $N_i=M_i\cap T$. Since $\mathfrak{m}=M_1\cap M_2$, $\mathfrak{m}$ is an ideal of $T$, and $N_1\cap N_2=(M_1\cap M_2)\cap T=\mathfrak{m}$. Moreover, $P_o=M_1\cap R_o\subseteq N_1$ and $Q_o=M_2\cap R_o\subseteq N_2$ since $R_o\subseteq T$. As $P_o$ and $Q_o$ are distinct maximal ideals of $R_o$, we have $P_o+Q_o=R_o$. Hence $1\in N_1+N_2$, and therefore, $N_1+N_2=T$. Consider the homomorphism $\varphi_i:R\longrightarrow T/N_i$ defined by $r\longmapsto r+N_i$. Since $\ker(\varphi_i)=N_i\cap R=\mathfrak{m}$, the homomorphism $\varphi_i$ induces an injective homomorphism $R/\mathfrak{m}\hookrightarrow T/N_i$. Thus, identifying $R/\mathfrak{m}$ with its image, each $T/N_i$ contains $R/\mathfrak{m}$. Since $P_o$ and $Q_o$ are comaximal, by the Chinese Remainder Theorem, the map $\rho:R_o\longrightarrow R_o/P_o\times R_o/Q_o$ is surjective. So, there exists $e\in R_o$ such that $\rho(e)=(0+P_o, 1+Q_o)$. Equivalently, $e\in P_o$ and $1-e\in Q_o$. Since $P_o\subseteq N_1$, the image of $e$ in $T/N_1$ is zero. On the other hand, $1-e\in Q_o\subseteq N_2$, so $(1-e)+N_2=0+N_2$, which is equivalent to $e+N_2=1+N_2$. Now, fix $t\in T$, and let $t_i=t+N_i\in T/N_i$ for $i=1,2$. If $t\in R$, then $t_i=t+N_i$ is the image of $t+\mathfrak{m}\in R/\mathfrak{m}$ under the embedding $R/\mathfrak{m}\hookrightarrow T/N_i$. Thus $t_i\in R/\mathfrak{m}$ for each $i$, and hence both $t_1$ and $t_2$ are algebraic over $R/\mathfrak{m}$. Suppose now that $t\notin R$. Since $R_o$ is the minimal overring of $R$, we have $R_o\subseteq R[t]$. Because $e\in R_o$, it follows that $e\in R[t]$. Hence there exist $r_0,\ldots,r_n\in R$ such that $e=r_0+r_1t+\cdots+r_nt^n$. Let $f(X)=r_0+r_1X+\cdots+r_nX^n\in R[X]$. Then, $e=f(t)$. If $\overline{r_i}=r_i+\mathfrak{m}\in R/\mathfrak{m}$ denotes the image of $r_i$ in $R/\mathfrak{m}$, define $\widetilde f(X)=\overline{r_0}+\overline{r_1}X+\cdots+\overline{r_n}X^n\in (R/\mathfrak{m})[X]$. Passing to $T/N_1$, the equality $e=f(t)$ becomes $e+N_1=\overline{r_0}+\overline{r_1}t_1+\cdots+\overline{r_n}t_1^n$. Since $e+N_1=0$, we obtain $\widetilde f(t_1)=0$. Likewise, passing to $T/N_2$, we obtain $e+N_2=\overline{r_0}+\overline{r_1}t_2+\cdots+\overline{r_n}t_2^n$. Since $e+N_2=1+N_2$, it follows that $\widetilde f(t_2)=1$. If $t_1$ were transcendental over $R/\mathfrak{m}$, since $\widetilde f(t_1)=0$, $\widetilde f=0$, but $\widetilde f(t_2)\neq 0$. Since $\widetilde f(t_2)=1$, define $g(X):=\widetilde f(X)-1\in (R/\mathfrak{m})[X]$. Then, $g(t_2)=0$. If $t_2$ were transcendental over $R/\mathfrak{m}$, this would force $g=0$, and hence $\widetilde f=1$. But then $\widetilde f(t_1)=1$, contradicting $\widetilde f(t_1)=0$, and hence both $t_1$ and $t_2$ are algebraic over $R/\mathfrak{m}$.		
		
		Choose monic polynomials $g_i\in(R/\mathfrak{m})[X]$ such that $g_i(t_i)=0$. Fix $i\in\{1,2\}$, and write
		\[
		g_i(X)=\overline{r_0}+\overline{r_1}X+\cdots+\overline{r_{n-1}}X^{n-1}+X^n.
		\]
		Choose representatives $r_j\in R$ such that $\overline{r_j}=r_j+\mathfrak{m}$, and put
		\[
		G_i(X)=r_0+r_1X+\cdots+r_{n-1}X^{n-1}+X^n \in R[X].
		\]
		To evaluate $g_i$ at $t_i\in T/N_i$, we use the embedding $ R/\mathfrak{m}\hookrightarrow T/N_i$ by $ r+\mathfrak{m}\longmapsto r+N_i.$ Thus, each coefficient $\overline{r_j}$ of $g_i$
		is viewed as $\varphi_i(r_j)=r_j+N_i$ in $T/N_i$. Consequently, $$G_i(t)+N_i =(r_0+N_i)+(r_1+N_i)t_i+\cdots +(r_{n-1}+N_i)t_i^{n-1}+t_i^n=g_i(t_i)=0. $$ Therefore, $G_i(t)\in N_i$ for each $i=1,2$. For $G=G_1G_2$, let $c:=G(t)\in N_1\cap N_2= \mathfrak{m} \subseteq R$, and define $h(X)=G(X)-c\in R[X]$. Then, $t$ is a root of the monic polynomial $h(X)$. Consequently, $T\subseteq\overline R$. The reverse inclusion follows from \cite[Theorem 9.3]{huckaba}. Hence $T=\overline R$.
		
		\smallskip We claim that $\overline R$ has exactly two regular maximal ideals, namely $N_1$ and $N_2$. Since $\overline R$ is integral over $R$ and $N_i\cap R=\mathfrak{m}$, both $N_i$ are maximal ideals. They are regular because they contain $\mathfrak{m}$, and they are distinct because their contractions to $R_o$ are $P_o$ and $Q_o$, respectively. Let $L$ be any regular maximal ideal of $\overline R$. By \cite[Theorem 2.4]{marot}, $L\cap R$ is regular. Integrality shows that $L\cap R$ is maximal, so $L\cap R=\mathfrak{m}$. Thus $N_1N_2=N_1\cap N_2=\mathfrak{m}\subseteq L$. Since $L$ is prime, either $N_1\subseteq L$ or $N_2\subseteq L$. Maximality gives $L=N_1$ or $L=N_2$. Therefore, $N_1,N_2$ are precisely the regular maximal ideals of $\overline R$.
		
		\smallskip We claim that $V_1$ and $V_2$ are the only valuation overrings of $R$ centered at $\mathfrak{m}$. Let $W$ be a Marot valuation overring of $R$ centered at $\mathfrak{m}$, with regular maximal ideal $M_W$. By \cite[Theorem 9.3]{huckaba} we have $R_o\subseteq \overline R\subseteq W$. The center of $W$ on $R_o$ is a regular maximal ideal, hence $M_W\cap R_o$ is equal to $P_o$ or $Q_o$. Without loss of generality, assume $M_W\cap R_o=P_o$. We know that the centers of the pair $(W,V_2)$ on $R_o$ are distinct. Therefore, since $M_1\cap M_2=\mathfrak{m}$ and $\overline{R} =V_1\cap V_2$, $M_W\cap M_2=\mathfrak{m}$ and $W\cap V_2=\overline R=V_1\cap V_2$. If $x\in\Reg(V_1\setminus W)$, then $x\notin V_2$ because otherwise $x\in\overline R\subseteq W$. So, $x^{-1}\in M_W\cap M_2=\mathfrak{m}\subseteq M_1$, contradicting to $x\in V_1$. Conversely, if $x\in\Reg(W\setminus V_1)$, then $x\notin V_2$, and $x^{-1}\in M_1\cap M_2=\mathfrak{m}\subseteq M_W$, contradicting to $x\in W$. Since $R$ is Marot, we have $W=V_1$.
		
		\smallskip We claim that $\overline{R}_{(N_i)}=V_i$. The Marot valuation overrings of $\overline R$ centered at $N_i$ consist only of $V_i$. Lemma \ref{centered} yields to $\overline{\overline R_{(N_i)}}=V_i$. Since regular localizations of an integrally closed ring are integrally closed, $\overline R_{(N_i)}=V_i$. It follows that $\overline R$ is Prüfer by \cite[Theorem 6.2]{huckaba}.
	\end{proof}

We now consider an integral comparable overring $R_o$ with two regular maximal ideals $P_o,Q_o$. We will show that $R_o$ is a minimal overring of $R_1=R+(P_o\cap Q_o)$ so that we can apply the preceding proposition to $R_1$.

\begin{proposition}\label{comparable-two}
		Let $R$ be a Marot ring with unique regular maximal ideal $\mathfrak{m}$, and let $R_o$ be a comparable overring with $R_o\subseteq\overline R$. Suppose $R_o$ has exactly two regular maximal ideals $P_o,Q_o$. Set $I=P_o\cap Q_o$ and $R_1=R+I$. Then:
		\begin{enumerate}[label=(\roman*)]
			\item $\overline R$ is a Marot Prüfer ring with exactly two regular maximal ideals $N_1,N_2$, and $N_1\cap N_2=I$.
			\item $R_1$ is the pullback $(\overline R;I;R/\mathfrak{m})$.
			\item $R_o$ is the minimal overring of $R_1$.
			\item If $R\subseteq S\subseteq\overline R$, then either $S\subsetneq R_o$ and $S$ has a unique regular maximal ideal, or $R_o\subseteq S$ and $S$ has exactly two regular maximal ideals.
			\item If $I\neq \mathfrak{m}$, then $R_1$ is a comparable overring of $R$.
		\end{enumerate}
	\end{proposition}
	\begin{proof}
		Since $R_o$ is integral over $R$, we have $P_o\cap R=Q_o\cap R=\mathfrak{m}$. Since $P_o$ and $Q_o$ are incomparable, $P_o$ is not $r$-divided by \cite[Theorem 2.8]{marot}. By \cite[Theorem 2.10]{marot}, $P_o\subsetneq P_o(R_o)_{(P_o)}$. Applying Lemma \ref{comparable-ideal} gives $R_o=R+P_o$. Similarly, $R_o=R+Q_o$.
		
		Since $I\cap R=\mathfrak{m}$ and $I$ is contained in every regular maximal ideal of $R_o$, Corollary \ref{regular-pullback} shows that $R_1$ has unique regular maximal ideal $I$ and that $P_o\cap R_1=Q_o\cap R_1=I$.
		
		\textbf{(iii)} Since $I\subseteq R_o$, we have $R_1=R+I\subseteq R_o$. If $R_1=R_o$, the equalities $P_o\cap R_1=Q_o\cap R_1=I$ would give $P_o=Q_o=I$, a contradiction. Thus, $R_1\subsetneq R_o$. Let $R_1\subsetneq S\subseteq R_o$. We first construct elements $a\in P_o\cap S$, and $b\in Q_o\cap S$ such that $a+b=1$. Choose $x\in S\setminus R_1$. Using $R_o=R+P_o$, write $x=r+p$ with $r\in R$ and $p\in P_o$. Since $p=x-r\in S$, we have $p\in P_o\cap S$. Moreover, $p\notin Q_o$; otherwise $p\in P_o\cap Q_o=I$, and hence $x=r+p\in R+I=R_1$, contrary to the choice of $x$. Using $R_o=R+Q_o$, write $p=u+q$ with $u\in R$ and $q\in Q_o$. Since $p\notin Q_o$, we have $u\notin Q_o\cap R=\mathfrak{m}$. Because $u+\mathfrak{m}$ is a non-zero element of the field $R/\mathfrak{m}$, there exists an element $v\in R$ such that $(u+\mathfrak{m})(v+\mathfrak{m})=1+\mathfrak{m}$. Therefore, $uv-1\in \mathfrak{m}$. Set $a:=vp\in P_o\cap S$ and $b:=1-a$. Since $a\in S$, we also have $b\in S$. Furthermore, $b=1-vp=1-v(u+q)=(1-uv)-vq$. Here $1-uv\in \mathfrak{m}\subseteq Q_o$, and $vq\in Q_o$. Therefore, $b\in Q_o\cap S$. Thus, we have $a\in P_o\cap S$ and $b\in Q_o\cap S$, with $a+b=1$. Now, for any $y\in R_o$, write $y=r_p+p'$ with $r_p\in R$, $p'\in P_o$. Then, $yb=r_pb+p'b\in S$ since $p'b\in P_oQ_o\subseteq I\subseteq S$. Similarly, writing $y=r_q+q'$ with $r_q\in R$, $q'\in Q_o$, we obtain $ya=r_qa+q'a\in S$. Hence, $y=ya+yb\in S$, so $S=R_o$. Therefore, $R_1\subsetneq R_o$ is minimal. Every overring of $R_1$ is also an overring of $R$, hence is comparable to $R_o$. Minimality now shows that every proper overring of $R_1$ contains $R_o$.
		
		\textbf{(i) \& (ii)} Since $R_1$ is integral over $R$, we have $\overline{R_1}=\overline R$. Apply Proposition \ref{minimal-two} to $R_1$ whose unique regular maximal ideal is $I$ and whose minimal overring is $R_o$. This proves (i) and shows that $I$ is an ideal of $\overline R$. After relabeling, we may assume $N_1\cap R_o=P_o$ and $N_2\cap R_o=Q_o$. Under the quotient map $\overline R\to\overline R/I$, the image of $R$ is isomorphic to $R/\mathfrak{m}$, and its preimage is exactly $R+I=R_1$. This proves (ii).
		
		\textbf{(iv)} Every such $S$ is comparable to $R_o$. If $S\subsetneq R_o$, then $R_o$ is an integral comparable overring of $S$. By \cite[Proposition 3.3]{marot}, $S$ has a unique regular maximal ideal. If $R_o\subseteq S\subseteq\overline R$, integrality of $\overline R$ over $S$ shows that the regular maximal ideals of $S$ are the contractions of $N_1,N_2$. These contractions are distinct since they contract further to $P_o,Q_o$ on $R_o$. Thus, $S$ has exactly two regular maximal ideals.
		
		\textbf{(v)} Let $T$ be any overring of $R$. If $R_o\subseteq T$, then $R_1\subseteq T$. Otherwise, $T\subsetneq R_o$. Suppose that $T\nsubseteq R_1$. Since $I$ is an ideal of $R_o$, the ring $T+I$ satisfies $R_1\subsetneq T+I\subseteq R_o.$ Part (iii) implies that $T+I=R_o$. By (iv), $T$ has a unique regular maximal ideal $\mathfrak n$. Since $R_o$ is integral over $T$, both $P_o$ and $Q_o$ contract to $\mathfrak n$. Choose $a\in P_o$ and $b\in Q_o$ with $a+b=1$, and write $a=t+i$, where $t\in T$ and $i\in I$. Then, $t=a-i\in P_o \cap T =\mathfrak n$ and $1-t=b+i\in Q_o\cap T=\mathfrak n.$ Thus $1\in\mathfrak n$, a contradiction. Therefore, $T\subseteq R_1$, and every overring of $R$ is comparable to $R_1$. If $I\neq\mathfrak m$, then $R_1\neq R$. Together with $R_1\subsetneq R_o\subsetneq T(R)$, this proves (v).
    \end{proof}

	The following theorem generalizes \cite[Theorem 24]{ayache}. As in Ayache's proof, we start with an arbitrary overring $T$ of $R$ and set $S=T\cap\overline R$. Ayache uses results on residually algebraic pairs of domains to show that $(S,T)$ is normal. These results do not apply directly to rings with zero divisors. Lemma \ref{normal-intersection} shows that $(S,T)$ is a normal pair by using $P$-extensions. If $S\subsetneq R_o$, normal closedness gives $T=S$. Otherwise, $R_o \subseteq S\subseteq T$. Thus conditions on the rings between $R$ and $\overline R$ allow us to determine whether every overring of $R$ is comparable to $R_o$.

	\begin{theorem}\label{comparable-two-characterization}
		Let $R$ be a Marot ring with unique regular maximal ideal $\mathfrak{m}$, and let $R_o$ be a proper overring of $R$ such that $R_o\subseteq\overline R$ and $R_o$ has exactly two regular maximal ideals $P_o,Q_o$. Then $R_o$ is a comparable overring of $R$ if and only if:
		\begin{enumerate}[label=(\roman*)]
			\item $\overline R$ is a Marot Prüfer ring;
			\item every intermediate ring $S$ between $R$ and $\overline R$ is comparable to $R_o$ under inclusion;
			\item every intermediate ring $S$ with $R\subseteq S\subsetneq R_o$ is normally closed.
		\end{enumerate}
	\end{theorem}
	\begin{proof}
		Suppose that $R_o$ is comparable. Proposition \ref{comparable-two} gives (i), and (ii) follows immediately from comparability. Since $R_o$ is integral over $R$, Lemma \ref{integral-comparable-normally-closed} gives (iii).
		
		Conversely, let $T$ be any overring of $R$. Set $S=T\cap\overline R$. By (ii), either $R_o\subseteq S$, in which case $R_o\subseteq T$, or $S\subsetneq R_o$. In the latter case, Lemma \ref{normal-intersection} shows that $(S,T)$ is normal, while (iii) says that $S$ is normally closed. Therefore, $T=S\subsetneq R_o$. Thus, every overring of $R$ is comparable to $R_o$.
	\end{proof}
	The following theorem generalizes \cite[Theorem 25]{ayache}. It describes when $R$ has a minimal overring with two regular maximal ideals. It also connects intermediate subfields of $\overline R/\mathfrak m$ with intermediate rings having a unique regular maximal ideal. To prove $(\mathrm{iii})\Rightarrow(\mathrm{i})$, we choose $e\in N_1$ with $1-e\in N_2$ and put $R_o=R[e]$. Since $e^2-e\in\mathfrak m$, $R_o$ is integral over $R$. We then show that every ring $T$ with $R\subsetneq T\subseteq\overline R$ contains $e$, and hence contains $R_o$. Corollary \ref{integral-minimal-characterization} shows that $R_o$ is the minimal overring of $R$. This avoids the separate localization argument used by Ayache for overrings outside $\overline R$.
	\begin{theorem}\label{minimal-two-characterization}
		Let $R$ be a Marot ring with unique regular maximal ideal $\mathfrak{m}$. The following are equivalent:
		\begin{enumerate}[label=(\roman*)]
			\item $R$ has a minimal overring $R_o$ with exactly two regular maximal ideals.
			\item $\overline R$ is a Marot Prüfer ring with exactly two regular maximal ideals $N_1,N_2$ such that $N_1\cap N_2=\mathfrak{m}$, and the image of $R/\mathfrak{m}$ in $\overline R/\mathfrak{m}$ is not properly contained in any subfield of $\overline R/\mathfrak{m}$.
			\item $\overline R$ is a Marot Prüfer ring with exactly two regular maximal ideals $N_1,N_2$such that $N_1\cap N_2=\mathfrak{m}$, and no intermediate ring $S$ with $R\subsetneq S\subseteq\overline R$ has a unique regular maximal ideal.
		\end{enumerate}
	\end{theorem}
	\begin{proof}
		\emph{(i)$\Rightarrow$(ii).}
		Proposition \ref{minimal-two} shows that $\overline R$ is a Marot Prüfer ring with exactly two regular maximal ideals $N_1,N_2$ and $N_1\cap N_2=\mathfrak{m}$. Let $\pi:\overline R\to\overline R/\mathfrak{m}$ be the canonical projection, and let $D$ be a subfield of $\overline R/\mathfrak{m}$ containing $R/\mathfrak{m}$. Set $S=\pi^{-1}(D)$. Since $S/\mathfrak{m}\cong D$, the ideal $\mathfrak{m}$ is a regular maximal ideal of $S$. The extension $S\subseteq\overline R$ is integral, so every regular maximal ideal of $S$ is the contraction of $N_1$ or $N_2$. Both contractions contain $\mathfrak{m}$ and hence equal $\mathfrak{m}$. Thus, $S$ has a unique regular maximal ideal.
		
		By Proposition \ref{comparable-two}(iv), $R_o\subseteq S$ would give two regular maximal ideals of $S$, a contradiction. So, comparability gives $S\subsetneq R_o$, and minimality yields to $S=R$. Consequently, $D=R/\mathfrak{m}$.
		
		\emph{(ii)$\Rightarrow$(iii).} Assume (ii), and let $R\subseteq S\subseteq\overline R$. Suppose that $S$ has a unique regular maximal ideal. Since $\overline R$ is integral over $S$, the regular maximal ideals of $S$ are the contractions $N_1\cap S$ and $N_2\cap S$. So, these contractions must be equal. Since $ (N_1\cap S)\cap(N_2\cap S) =(N_1\cap N_2)\cap S=\mathfrak{m}, $ both contractions equal $\mathfrak{m}$. Thus, $S/\mathfrak{m}$ is a subfield of $\overline R/\mathfrak{m}$ containing $R/\mathfrak{m}$. By (ii), we have $S/\mathfrak{m}=R/\mathfrak{m}$. Hence, $S=R$.
		
		\emph{(iii)$\Rightarrow$(i).}
		Since $N_1,N_2$ are distinct maximal ideals, $N_1+N_2=\overline R$. Choose $e\in N_1$ such that $1-e\in N_2$. Then
		\[
		e^2-e\in N_1\cap N_2=\mathfrak{m}\subseteq R.
		\]
		Moreover, $e\notin R$: otherwise $e\in N_1\cap R=\mathfrak{m}\subseteq N_2$, contradicting to $1-e\in N_2$. Set $R_o=R[e]$. This is a proper integral overring of $R$ contained in $\overline R$. Its regular maximal ideals are the contractions $P_o=N_1\cap R_o$ and $Q_o=N_2\cap R_o$. They are distinct because $e\in P_o$ and $1-e\in Q_o$. Hence, $R_o$ has exactly two regular maximal ideals.
		
		We first prove that every intermediate ring $T$ with $R\subsetneq T\subseteq\overline R$ contains $R_o$. Let $R\subsetneq T\subseteq\overline R$. The regular maximal ideals of $T$ are the contractions of $N_1,N_2$. By (iii), these contractions are distinct. Choose $a\in N_1\cap T$ and $b\in N_2\cap T$ with $a+b=1$. Since
		\[
		e-a\in N_1,\qquad e-a=b-(1-e)\in N_2,
		\]
		we have $e-a\in\mathfrak{m}\subseteq T$. Thus $e\in T$, and $R_o=R[e]\subseteq T$. Thus every intermediate ring $T$ with $R\subsetneq T\subseteq\overline R$ contains $R_o$. Since $\overline R$ is Prüfer, Corollary \ref{integral-minimal-characterization} shows that $R_o$ is the minimal overring of $R$.
    \end{proof}

The next consequences describe the position of $R_o$ and $R_1$ among integral comparable overrings and characterize when $R_1=R$.

\begin{corollary}\label{largest-integral-comparable}
		Let $R$ be a Marot ring and let $R_o\subseteq\overline R$ be a comparable overring having exactly two regular maximal ideals. Then:
		\begin{enumerate}[label=(\roman*)]
			\item Every comparable overring of $R$ that is integral over $R$ is contained in $R_o$. Thus $R_o$ is the largest integral comparable overring of $R$.
			\item $R_o$ is the unique comparable overring of $R$ having exactly two regular maximal ideals.
		\end{enumerate}
	\end{corollary}
	\begin{proof}
		\textbf{(i)} Let $S$ be an integral comparable overring of $R$. Comparability implies that $S\subseteq R_o$ or $R_o\subsetneq S$. In the latter case, $S$ is an integral comparable overring of $R_o$. By \cite[Proposition 3.3]{marot}, $R_o$ would have a unique regular maximal ideal, a contradiction.
		
		\textbf{(ii)} Let $T$ be a comparable overring of $R$ having exactly two regular maximal ideals. It is not a Marot valuation ring, so \cite[Proposition 3.2]{marot} implies that it is integral over $R$. By (i), $T\subseteq R_o$. If $T\subsetneq R_o$, then $R_o$ is an integral comparable overring of $T$. Again, \cite[Proposition 3.3]{marot} would give a unique regular maximal ideal of $T$, a contradiction. Thus, $T=R_o$.
	\end{proof}

\begin{corollary}\label{two-minimal-equivalences}
		Let $R$ be a Marot ring with unique regular maximal ideal $\mathfrak{m}$, and let $R_o\subseteq\overline R$ be a comparable overring having exactly two regular maximal ideals $P_o,Q_o$. Set $I=P_o\cap Q_o$ and $R_1=R+I$. The following conditions are equivalent:
		\begin{enumerate}[label=(\roman*)]
			\item $R_o$ is the minimal overring of $R$.
			\item $I=\mathfrak{m}$.
			\item $R_1=R$.
			\item $R_o$ is the only integral comparable overring of $R$.
			\item $R_o$ is the unique comparable overring of $R$ that is not a valuation ring.
		\end{enumerate}
	\end{corollary}
	\begin{proof}
		(i)$\Rightarrow$(ii). Since $P_o\cap R=Q_o\cap R=\mathfrak{m}$, Lemma \ref{minimal-ideals} gives $I=P_o\cap Q_o=\mathfrak{m}$.
		
		(ii)$\Rightarrow$(iii). We have $R_1=R+I=R+\mathfrak{m}=R$.
		
		(iii)$\Rightarrow$(i). By Proposition \ref{comparable-two}(iii), $R_o$ is the minimal overring of $R_1$. Since $R_1=R$, the assertion follows.
		
		(i)$\Rightarrow$(iv). Let $S$ be an integral comparable overring of $R$. By Corollary \ref{largest-integral-comparable}(i), $S\subseteq R_o$. Since $S$ is a proper overring of $R$ and $R_o$ is the minimal overring, $R_o\subseteq S$. Hence, $S=R_o$.
		
		(iv)$\Rightarrow$(v). The ring $R_o$ is not a Marot valuation ring since it has two regular maximal ideals. Let $S$ be a comparable overring of $R$ that is not a valuation ring. By \cite[Proposition 3.2]{marot}, $S$ is integral over $R$. Thus, (iv) implies that $S=R_o$.
		
		(v)$\Rightarrow$(ii). Suppose that $I\neq\mathfrak{m}$. By Proposition \ref{comparable-two}, $R_1$ is a comparable overring of $R$ satisfying $R_1\subsetneq R_o\subseteq\overline R$ and $\overline{R_1}=\overline R$. If $R_1$ were a Marot valuation ring, it would be integrally closed, so $R_1=\overline{R_1}=\overline R$, a contradiction. Thus, $R_1$ is a comparable overring different from $R_o$ which is not a Marot valuation ring, contradicting (v). Therefore, $I=\mathfrak{m}$.
	\end{proof}

	\begin{corollary}\label{largest-one-maximal}
		Let $R$ be a Marot ring with unique regular maximal ideal $\mathfrak{m}$, and let $R_o\subseteq\overline R$ be a comparable overring having exactly two regular maximal ideals $P_o,Q_o$. Set $I=P_o\cap Q_o$ and $R_1=R+I$, and suppose that $I\neq\mathfrak{m}$. Then $R_1$ is the largest integral comparable overring of $R$ different from $R_o$. Equivalently, $R_1$ is the largest integral comparable overring of $R$ having a unique regular maximal ideal.
	\end{corollary}
	\begin{proof}
		By Proposition \ref{comparable-two}, $R_1$ is an integral comparable overring of $R$ with unique regular maximal ideal $I$ satisfying $R_1\subsetneq R_o$. Moreover, $R_o$ is the minimal overring of $R_1$. Let $S\neq R_o$ be an integral comparable overring of $R$. Corollary \ref{largest-integral-comparable}(i) implies that $S\subsetneq R_o$. Since $R_1$ is comparable, either $S\subseteq R_1$ or $R_1\subsetneq S$. In the latter case, the minimal overring property implies that $R_o\subseteq S$, contradicting to $S\subsetneq R_o$. Thus, $S\subseteq R_1$, proving the first assertion. By Proposition \ref{comparable-two}(iv), every integral comparable overring properly contained in $R_o$ has a unique regular maximal ideal whereas $R_o$ has two.
	\end{proof}

\begin{corollary}\label{two-closure-comparable}
		Let $R$ be a Marot ring, and let $R_o\subseteq\overline R$ be a comparable overring having exactly two regular maximal ideals. Then $\overline R$ is a comparable overring of $R$ if and only if $R_o=\overline R$.
	\end{corollary}
	\begin{proof}
		Suppose that $\overline R$ is a comparable overring of $R$. Since it is integral over $R$, Corollary \ref{largest-integral-comparable}(i) implies that $\overline R\subseteq R_o$. Together with $R_o\subseteq\overline R$, this yields $R_o=\overline R$. The converse follows immediately from the comparability of $R_o$.
	\end{proof}

\begin{corollary}\label{two-minimal-other-comparables}
		Let $R$ be a Marot ring with unique regular maximal ideal $\mathfrak{m}$, and let $R_o$ be its minimal overring. Suppose that $R_o$ has exactly two regular maximal ideals. Then every comparable overring $S\neq R_o$ of $R$ is a Marot valuation ring satisfying $\overline R\subsetneq S$.
	\end{corollary}
	\begin{proof}
        Since $R_o$ has two regular maximal ideals, it is not a Marot valuation ring. Hence, \cite[Proposition 3.2]{marot} implies that $R_o\subseteq\overline R$. By Proposition \ref{minimal-two}, $\overline R$ has exactly two regular maximal ideals. Corollary \ref{two-minimal-equivalences} shows that $S$ is a Marot valuation ring. Since $S$ is integrally closed containing $R$, we have $\overline R\subseteq S$. $R$ is proper in $S$ since $S$ has a unique regular maximal ideal. Therefore, $\overline R\subsetneq S$.
    \end{proof}

We next compute $(I:I)$, $(R_1:I)$, and the conductors of $R_1$ in $R_o$ and $\overline R$.

\begin{proposition}\label{two-colon}
		Let $R$ be a Marot ring with unique regular maximal ideal $\mathfrak{m}$, and let $R_o\subseteq\overline R$ be a comparable overring having exactly two regular maximal
		ideals $P_o,Q_o$. Set $I=P_o\cap Q_o$ and $R_1=R+I$. Then
		\[
		\overline R=(I:I)=(R_1:I),\qquad
		(R_1:R_o)=(R_1:\overline R)=I.
		\]
		In particular, $I$ is a divisorial ideal of $R_1$: $ I_v:=(R_1:(R_1:I))=I$. \\ If $R_o$ is the minimal overring of $R$, then
		$I=\mathfrak{m}$, $R_1=R$, and consequently
		\[
		\overline R=(\mathfrak{m}:\mathfrak{m})=(R:\mathfrak{m}),
		\qquad (R:\overline R)=\mathfrak{m},
		\qquad \mathfrak{m}_v=\mathfrak{m}.
		\]
	\end{proposition}
	
	\begin{proof}
		By Proposition \ref{comparable-two}, the ring $R_1=R+I$ has a unique regular maximal ideal $I$, its minimal overring is $R_o$, and $\overline{R_1}=\overline R$. Moreover, $I$ is an ideal of $\overline R$, and $\overline R$ has exactly two regular maximal ideals $N_1,N_2$, with $N_1\cap N_2=I$.
		
		Applying Proposition \ref{minimal-two} to $R_1$, we obtain Marot valuation overrings $V_1,V_2$, with regular maximal ideals $M_1,M_2$, such that $\overline R=V_1\cap V_2$, $V_i=(\overline R)_{(N_i)}$, and $M_i\cap\overline R=N_i$. Moreover, one can easily show that $N_iV_i=M_i$.
		
		Since $N_1,N_2$ are distinct maximal ideals of $\overline R$, we have $I=N_1\cap N_2=N_1N_2$. Let $s\in\Reg(N_1\setminus N_2)$. In the regular localization $V_2=(\overline R)_{(N_2)}$, the element $s$ is a unit. Since $s\in N_1$, this gives $1=ss^{-1}\in N_1V_2$. Hence, $N_1V_2=V_2$, and therefore, $$IV_2=(N_1N_2)V_2=(N_2V_2)(N_1V_2) =M_2V_2 =M_2.$$ Similarly, $IV_1=M_1$.
		
		\smallskip
		Now, we claim that $(I:I)=\overline R$. Since $I$ is an ideal of $\overline R$, we have $\overline R\subseteq(I:I)$. For the reverse inclusion, first let $x\in\Reg((I:I))$. Then $xI\subseteq I$. Observe that $xM_i=x(IV_i)=(xI)V_i\subseteq IV_i=M_i$. Suppose that $x\notin V_i$. Since $V_i$ is a Marot valuation ring, we have $x^{-1}\in M_i$. Since $xM_i\subseteq M_i$, we have $1=xx^{-1}\in M_i$, which is impossible. Therefore, $x\in V_i$. So, $x\in V_1\cap V_2=\overline R$. Since $R$ is Marot, we have $(I:I)\subseteq\overline R$. Hence, $(I:I)=\overline R$.
		
		\smallskip
		Now, we claim that $(R_1:I)=\overline R$. We have $\overline R\subseteq(R_1:I)$ because $\overline RI\subseteq I\subseteq R_1$. Let $x\in\Reg((R_1:I))$. Then, $xI\subseteq R_1$, and $r(xa)=x(ra)\in xI$ for every $r\in R_1$, and $a\in I$. Thus, $xI$ is a regular ideal of $R_1$. If $xI\nsubseteq I$, then $xI=R_1$ since every proper regular ideal of $R_1$ is contained in $I$. Hence, $I=x^{-1}R_1$, and therefore, $(I:I)=(x^{-1}R_1:x^{-1}R_1)=(R_1:R_1)=R_1.$ This contradicts to $(I:I)=\overline R$ so that $R_1\subsetneq R_o\subseteq\overline R$. Thus $xI\subseteq I$, so $x\in(I:I)=\overline R$. Since $R$ is Marot, we have $(R_1:I)\subseteq\overline R$. Hence, $(R_1:I)=\overline R$.
		
		\smallskip
		Now, we claim that $(R_1:R_o)=I$. If $x\in(R_1:R_o)$, then $xR_o\subseteq R_1$. Since $1\in R_o$, we have $x=x\cdot1\in R_1$. Thus, $(R_1:R_o)\subseteq R_1$. Moreover, $(R_1:R_o)$ is an ideal of $R_1$. Since $I$ is an ideal of $\overline R$ and $R_o\subseteq\overline R$, we have $IR_o\subseteq I\subseteq R_1$. Therefore $I\subseteq(R_1:R_o)$. The ideal $(R_1:R_o)$ is proper. Otherwise, $1\in(R_1:R_o)$ which would imply that $R_o=1R_o\subseteq R_1$, contradicting to $R_1\subsetneq R_o$. Maximality of $I$ forces $(R_1:R_o)=I$. Moreover, we have $I\subseteq(R_1:\overline R)\subseteq(R_1:R_o)=I$. Hence, $(R_1:\overline R)=I$.
		
		\smallskip
		Now, we claim that $I$ is a divisorial ideal of $R_1$. By using the equalities proved above, we obtain $I_v=(R_1:(R_1:I))=(R_1:\overline R)=I$. Hence, $I$ is a divisorial ideal of $R_1$.
		
		\smallskip
		Finally, suppose that $R_o$ is the minimal overring of $R$. By Corollary \ref{two-minimal-equivalences}, the result follows.
	\end{proof}

\subsection{Type III}\label{subsec-type-iii}
We now consider the case in which $R_o$ has a unique regular maximal ideal $Q_o$ properly containing $\mathfrak m$. The conductor criterion in the following proposition holds whenever $R_o$ has a unique regular maximal ideal, so it also covers Type I. It gives a generalization of \cite[Proposition 27]{ayache} for Marot rings. We do not obtain a general characterization of comparable overrings of Type III.

\begin{proposition}\label{radical-conductor-comparable}
		Let $R$ be a Marot ring with unique regular maximal ideal $\mathfrak{m}$, and let $R_o$ be a proper overring of $R$ such that $R_o\subseteq\overline R$. Suppose that $R_o$ has a unique regular maximal ideal $Q_o$. If
		\begin{enumerate}[label=(\roman*)]
			\item $\overline R$ is Prüfer;
			\item every intermediate ring $S$ between $R$ and $\overline R$ is comparable to $R_o$ under inclusion;
			\item the radical of the conductor $(R:R_o)$ in $R_o$ is its unique regular maximal ideal $Q_o$; explicitly, \[ Q_o=\sqrt{(R:R_o)} =\{x\in R_o:x^n\in(R:R_o) \text{ for some integer }n\geq1\}, \]
		\end{enumerate}
		then $R_o$ is a comparable overring of $R$.
	\end{proposition}
	\begin{proof}
		Let $T$ be any overring of $R$, and put $S=T\cap\overline R$. By (ii), either $R_o\subseteq S$, in which case $R_o\subseteq T$, or $S\subsetneq R_o$. Suppose that $S\subsetneq R_o$. Since $S\subseteq R_o$ is integral, $S$ has a unique regular maximal ideal $\mathfrak{n}=Q_o\cap S.$ By (i) and Lemma \ref{normal-intersection}, the pair $(S,T)$ is normal. Suppose that $T\nsubseteq R_o$. So, there exists an element $t\in\Reg(T\setminus R_o)$. In particular, $t\notin S$. By Lemma \ref{normal-regular-inverse}, $t^{-1}\in\mathfrak{n}=Q_o\cap S\subseteq Q_o.$ So, part (iii) implies that there is a positive integer $n$ such that $(t^{-1})^n\in(R:R_o). $ Hence, $(t^{-1})^nx\in R$ for every $x\in R_o$, and consequently $x=t^n((t^{-1})^nx)\in R[t]\subseteq T.$ Thus, $R_o\subseteq T$. Since $R_o\subseteq\overline R$, this implies $R_o\subseteq T\cap\overline R=S, $ contradicting to $S\subsetneq R_o$. Therefore, $T\subseteq R_o$. Thus, every overring of $R$ is comparable to $R_o$.
	\end{proof}
	Next, we give an example which shows the converse of Proposition \ref{radical-conductor-comparable} does not hold, in general, even when the conditions $(i)$ and $(ii)$ hold. Throughout the example, derivatives are understood algebraically. Polynomial differentiation is defined by $(X^n)'=nX^{n-1}$ and extended to rational functions by
	\[\left(\frac pq\right)'=\frac{p'q-pq'}{q^2}.\]
	No limits or convergence are involved.
	\begin{example}\label{radical-conductor-converse}
		Consider
		\[\mathbb Z_{(2)}=\left\{\frac ab\in\mathbb Q:a,b\in\mathbb Z,\ b\text{ is odd}\right\}
		\hspace{1mm} \text{ and }\hspace{3mm} W=\mathbb Q[X]_{(X)}=\left\{\frac{p(X)}{q(X)}:p,q\in\mathbb Q[X],\ q(0)\neq0\right\}.\]
		Define
		\[ R=\{f\in W:f(0)\in \mathbb Z_{(2)},\ f'(0)\in2\mathbb Z_{(2)}\}, \qquad R_o=\{f\in W:f(0)\in \mathbb Z_{(2)}\}. \]
		Clearly, $R$ is a subring of $R_o$. Both rings have the quotient field $\mathbb Q(X)$. Moreover, we have $X\in R_o\setminus R$ and $X^{-1}\notin R_o,$ so $R\subsetneq R_o\subsetneq\mathbb Q(X).$
		
		We first show that $R$ is local with maximal ideal $\mathfrak m =\{f\in R:f(0)\in2\mathbb Z_{(2)}\}. $ Consider the homomorphism $\varphi:R\longrightarrow \mathbb Z_{(2)}/2\mathbb Z_{(2)}$ defined by $ f\longmapsto f(0)+2\mathbb Z_{(2)}.$ Now we want to show that $\varphi$ is surjective. Moreover, every $a\in\mathbb Z_{(2)}$ belongs to $R$: its value at $0$ is $a$, and its derivative is zero. Thus, $\varphi(a)=a+2\mathbb Z_{(2)}, $ so $\varphi$ is surjective. An element $f\in R$ belongs to $\ker\varphi$ precisely when $f(0)\in2\mathbb Z_{(2)}$. Therefore, $\ker\varphi=\mathfrak m$, and the First Isomorphism Theorem implies that $R/\mathfrak m \cong\mathbb Z_{(2)}/2\mathbb Z_{(2)} \cong\mathbb F_2.$ Hence, $\mathfrak m$ is a maximal ideal of $R$.
		
		If $f\in R\setminus\mathfrak m$, then $f(0)$ is a unit of $\mathbb Z_{(2)}$. Consequently, $\displaystyle\frac{1}{f}:=f^{-1}\in W$ and 
		\[(f^{-1})(0)=f(0)^{-1}\in \mathbb Z_{(2)},\qquad (f^{-1})'(0) =-\frac{f'(0)} {f(0)^2}\in2\mathbb Z_{(2)}. \]
		Thus, $f^{-1}\in R$. Therefore, $\mathfrak m$ is the unique maximal ideal of $R$. In particular, $R$ is a local domain, and hence a Marot ring. On the other hand, the same argument, without the derivative condition, shows that $R_o$ is local with maximal ideal \[ Q_o=\{f\in R_o:f(0)\in2\mathbb Z_{(2)}\}. \]
		
		We next show that $R_o$ is a valuation domain. Let $f\in\mathbb Q(X)\setminus R_o$. If $f\notin W$, then $f^{-1}\in W$ and $(f^{-1})(0)=0$, so
		$f^{-1}\in R_o$. Otherwise, $f\in W$ and $c=f(0)\notin \mathbb Z_{(2)}$.
		Then $$(f^{-1})(0) =\displaystyle\frac{1}{f(0)} =\displaystyle\frac{1}{c} =c^{-1}\in2\mathbb Z_{(2)},$$ so, $f^{-1}\in R_o$. Thus, $R_o$ is a valuation domain. Hence $R_o$ is integrally closed and Prüfer.
		
		Next, we show that $R_o$ is integral over $R$. Let $f\in R_o$, and put $a=f(0)\in\mathbb Z_{(2)}$. We observe that $a\in R$. Set $b:=(f-a)^2\in W.$ We show that $b\in R$. First, $b(0)=(f(0)-a)^2=0\in\mathbb Z_{(2)}.$ Moreover, differention implies that $b'=2(f-a)f',$ and hence $b'(0)=2(f(0)-a)f'(0)=0\in2\mathbb Z_{(2)}.$ Thus, $b\in R$. Since $a,b\in R$, the polynomial $P(Z):=(Z-a)^2-b =Z^2-2aZ+(a^2-b)$ belongs to $R[Z]$ and is monic. Hence, $P(f)=(f-a)^2-b=0.$ Therefore, $R_o$ is integral over $R$.
		
		Since $R_o$ is integral over $R$ and integrally closed, we conclude that $R_o=\overline R$. In particular, $\overline R$ is Prüfer, and every intermediate ring between $R$ and $\overline R$
		is comparable to $R_o$.\bigskip
		
		We now show that $R_o$ is comparable to every overring of $R$. Let $T$ be an overring of $R$ such that $T\nsubseteq R_o$, and choose $t\in T\setminus R_o$. Fix $g\in R_o$, and put $a=g(0)\in \mathbb Z_{(2)}$.
		
		First, suppose that $t\notin W$. Then $t^{-1}\in XW$, while $g-a\in XW$. Consequently, $$h:=\displaystyle\frac{g-a}{t}\in X^2W\subseteq R.$$ It follows that $g=a+th\in T. $ Next, suppose that $t\in W$. Then, $c=t(0)\notin \mathbb Z_{(2)}$, so $c^{-1}\in2\mathbb Z_{(2)}$. Write $c^{-1}=2u$ with $u\in\mathbb Z_{(2)}$. Since $g'(0)\in\mathbb Q$, we can write $g'(0)=\displaystyle\frac{v}{2^r}, $ where $v\in\mathbb Z_{(2)}$ and $r$ is a nonnegative integer. Taking $n=r+1\geq1$, we obtain 
		\[ \frac{g'(0)}{c^n} =\frac{v}{2^r}(2u)^{r+1} =2vu^{r+1}\in2\mathbb Z_{(2)}. \]
		
		Set $h:=\displaystyle\frac{g-a}{t^n}\in W.$ Since $g(0)-a=0$, we have $h(0)=0$, and $ h'(0)=\displaystyle\frac{g'(0)}{c^n}\in2\mathbb Z_{(2)}.$ Thus, $h\in R$, and therefore, $g=a+t^nh\in T.$ In both cases, $g\in T$. Hence, $R_o\subseteq T$, proving that $R_o$ is a comparable overring of $R$.
		
		\medskip Finally, we compute the conductor $(R:R_o)=\{f\in R:f(0)=0\}.$ For the first inclusion, let $f\in(R:R_o)$.
		Since $1\in R_o$, we have $f\in R$.
		Suppose that $f(0)\neq0$, and choose $N\geq1$ such that $2^{-N}f(0)\notin2\mathbb Z_{(2)}.$ Although $X/2^N\in R_o$, we have
		\[ \left(f\frac{X}{2^N}\right)'(0)=2^{-N}f(0)\notin2\mathbb Z_{(2)}. \]
		Thus $fX/2^N\notin R$, a contradiction. Therefore $f(0)=0$. Conversely, let $f\in R$ with $f(0)=0$. For every $g\in R_o$, we have $(fg)(0)=0$ and $ (fg)'(0)=f'(0)g(0)\in2\mathbb Z_{(2)}.$ Hence, $fg\in R$.\medskip
		
		If $f\in R_o$ and $f(0)=0$, then $f^2\in \{f\in R:f(0)=0\}=(R:R_o)$. If $f(0)\neq0$, no positive power of $f$ belongs to $(R:R_o)$. Consequently, taking the radical in $R_o$, we obtain
		\[
		\sqrt{(R:R_o)}
		=\{f\in R_o:f(0)=0\}
		\subsetneq Q_o.
		\]
		Indeed, $2\in Q_o$, but no positive power of $2$ belongs to $(R:R_o)$. Thus, $R_o$ is comparable and satisfies conditions (i) and (ii) of Proposition \ref{radical-conductor-comparable} while condition (iii) fails.
	\end{example}
	\section{Counterexamples for the proposed conjectures in \cite{ayache}} \label{counterexamples}
	
	In \cite[Problems 29 and 30]{ayache}, two questions are raised about integral comparable overrings. The following examples give negative answers to both questions.
	
	The first question, \cite[Problem 29]{ayache}, asks whether ``normally closed'' can be replaced by ``quasilocal'' in \cite[Theorem 24(iii)]{ayache}. In other words, under the other assumptions of that theorem, is it enough that every intermediate ring $S$ with $R\subseteq S\subsetneq R_o$ has a unique maximal ideal? The following example shows that this is not enough. Here $R_o=\overline R$ is Prüfer with exactly two maximal ideals, and every intermediate ring properly below $R_o$ is local, but $R_o$ is not comparable.
	\begin{example}\label{counterexample-problem-29}
		Put
		\[W= \left\{ \frac{p(X)}{q(X)}: p,q\in\mathbb Q[X],\ q(1)\neq0,\ q(0)\neq0 \right\},\]
		and define 
		\[ R=\{f\in W:f(0)=f(1)\in\mathbb Z_{(2)}\} \hspace{2mm}\text{ and }\hspace{2mm} R_o= \{f\in W:f(0),f(1)\in\mathbb Z_{(2)}\}. \]
		We show that $\overline R=R_o$ is Prüfer, that $R_o$ has exactly two maximal ideals and that overring $S$ of $R$ such that$R\subseteq S\subsetneq R_o$ is local. Nevertheless, $R_o$ is not a comparable overring of $R$. \looseness=-1
	
		Observe that \(R\subseteq R_o\subseteq W\). So,  $ T(R) \subseteq T(R_o) \subseteq T(W)$. Set $J=X(X-1)W$, a nonzero principal ideal of $W$.	Every $h\in J$ satisfies $h(0)=h(1)=0 \in\mathbb Z_{(2)} $, so $J\subseteq R$. Thus, for every $f\in W$,
		\[	f=\frac{X(X-1)f}{X(X-1)} \in T(R). \]
		So, $ T(W)\subseteq T(R).$ Consequently, $R$, $R_o$ and $W$ have the same quotient field $\mathbb Q(X)$. 
	
		Now, consider the surjective homomorphism $\rho:R_o\longrightarrow \mathbb Z_{(2)} \times \mathbb Z_{(2)}$ defined by \mbox{$f\longmapsto(f(0),f(1))$}. So, $\rho$ is onto. For any $a,b\in\mathbb Z_{(2)}$, the polynomial $a+(b-a)X$ has image $(a,b)$. Moreover, $\ker\rho=\{f\in R_o:f(0)=f(1)=0\}=J$.
	
		We first observe that $J\subseteq\operatorname{Jac}(R_o)$. Indeed, let $j\in J$ and $f\in R_o$. Write $1-jf=p/q$, where $p,q\in\mathbb Q[X]$ such that $q(0)\neq0$ and $q(1)\neq0$. Since $j(0)=j(1)=0$, evaluating at these points gives 
		\[ \frac{p(0)}{q(0)}=1, \qquad \frac{p(1)}{q(1)}=1. \]
		Therefore $p(0)\neq0$ and $p(1)\neq0$. Thus, $(1-jf)^{-1}=q/p\in W$. Its values at both points are $1\in\mathbb Z_{(2)}$, and hence $(1-jf)^{-1}\in R$. Since $J$ is contained in every maximal ideal of $R_o$, $R_o/J$ has the same number of maximal ideals as $R_o$. Since $R_o/J\cong\mathbb Z_{(2)}\times\mathbb Z_{(2)}$ has exactly two maximal ideals, $R_o$ also has exactly two maximal ideals:
		\[
		P_o=\{f\in R_o:f(0)\in2\mathbb Z_{(2)}\}\hspace{2mm}\text{ and }\hspace{2mm} Q_o=\{f\in R_o:f(1)\in2\mathbb Z_{(2)}\}.
		\]
		These are the inverse images under $\rho$ of $2\mathbb Z_{(2)}\times\mathbb Z_{(2)}$ and $\mathbb Z_{(2)}\times2\mathbb Z_{(2)}$, the two maximal ideals of $\mathbb Z_{(2)}\times\mathbb Z_{(2)}$, respectively.
	
		We next show that $R_o=\overline R$ is Prüfer. Consider the polynomial $g(X)=X$. Since $g(1)$ and $g(0)$ are elements of $\mathbb{Z}_{(2)}$, we have $X=g(X)\in R_o$. So, $R[X]\subseteq R_o$. Conversely, let $f\in R_o$, and put $a=f(0)$ and $b=f(1)$. Then, $a,b\in\mathbb Z_{(2)}\subseteq R$. The function $r:=f-a-(b-a)X$ satisfies $r(0)=r(1)=0$, so $r\in J\subseteq R$. Therefore $f=(r+a)+(b-a)X\in R[X]$ since both coefficients belong to $R$. Thus, $R_o\subseteq R[X]$, proving that $R_o=R[X]$. 
	
		Put $c:=X^2-X\in J\subseteq R$. The monic polynomial $P(Z)=Z^2-Z-c\in R[Z]$ satisfies $P(X)=X^2-X-c=0$. Thus, $X$ is integral over $R$, and hence $R_o=R[X]$ is integral over $R$. The localizations at the two maximal ideals are
		\[(R_o)_{P_o} = \{f\in\mathbb Q(X): f\text{ is defined at }0,\ 	f(0)\in\mathbb Z_{(2)}\},\]
		\[(R_o)_{Q_o} = \{f\in\mathbb Q(X): f\text{ is defined at }1,\  f(1)\in\mathbb Z_{(2)}\}. \]
		Let $f\in(R_o)_{P_o}$. Write $f=g/s$, where	$g,s\in R_o$ and $s\notin P_o$. Then, $g(0)\in\mathbb Z_{(2)}$ and $s(0)\in\mathbb Z_{(2)}\setminus2\mathbb Z_{(2)}$.	Thus, $s(0)$ is a unit of $\mathbb Z_{(2)}$, and hence $f(0)=\frac{g(0)}{s(0)}\in\mathbb Z_{(2)}.$ Conversely, let $f\in\mathbb Q(X)$ be defined at $0$ with $f(0)\in\mathbb Z_{(2)}$. Although $f$ is defined at $0$, it need not be defined at $1$. To obtain an element of $W$, we need a rational function defined at both points. We multiply by a power of $1-X$ since it equals $0$ at $X=1$ and $1$ at $X=0$. Since $f$ is defined at $0$, we can write $f=p/q$, where $p,q\in\mathbb Q[X]$ and $q(0)\neq0$.
	
		If $q(1)\neq0$, we let $m=0$ and $u=q$. If $q(1)=0$, then we factor out all factors $1-X$ from $q$ and write
		\[ q(X)=(1-X)^m u(X), \qquad m\geq1,\quad u(1)\neq0. \]
		Thus, in both cases, $q(X)=(1-X)^m u(X)$ with $m\geq0$ and $u(1)\neq0$.	Evaluating at $X=0$ implies that
		\[ q(0)=(1-0)^m u(0)=u(0). \]
		Since $q(0)\neq0$, we also have $u(0)\neq0$. Taking $n=m+1$, we obtain
		\[ (1-X)^nf=\frac{(1-X)p(X)}{u(X)}\in W. \]
		The denominator is nonzero at both points, and $1-X$ vanishes at $1$. Moreover, $\bigl((1-X)^nf\bigr)(0) =f(0)\in\mathbb Z_{(2)}.$  Therefore, $(1-X)^nf\in R_o$. Since $(1-X)^n$ has values $1$ and $0$ at $0$ and $1$, respectively, it belongs to $R_o\setminus P_o$.	Consequently,
		\[ f=\frac{(1-X)^nf}{(1-X)^n}\in(R_o)_{P_o}. \]
		The second equality is proved similarly.
	
		Both localizations are valuation domains. Indeed, let	$0\neq f\in\mathbb Q(X)\setminus(R_o)_{P_o}$. If it has a pole at $0$, its inverse has value zero there. Otherwise, $f$ is defined at $0$ and $c:=f(0)\in\mathbb Q\setminus\mathbb Z_{(2)}$. Then $1/f(0)=c^{-1}\in2\mathbb Z_{(2)}\subseteq\mathbb Z_{(2)},$ so $f^{-1}\in(R_o)_{P_o}$. Therefore, $(R_o)_{P_o}$ is a valuation domain. Similarly, $(R_o)_{Q_o}$ is a valuation domain. Therefore, $R_o$ is Prüfer and integrally closed. Since $R_o$ is integral over $R$, $R_o=\overline R$.
	
		Now, let $R\subseteq S\subsetneq R_o$, and define $\mathfrak a= \{f(1)-f(0):f\in S\}. $ Since $S$ contains all constant functions in $\mathbb Z_{(2)}$, the set $\mathfrak a$ is an ideal of $\mathbb Z_{(2)}$. We claim that $S=\{f\in R_o:f(1)-f(0)\in\mathfrak a\}.$ We observe that $f(1)-f(0)\in\mathfrak a$ for every $f\in S$. Conversely, let $f\in R_o$ such that $f(1) -f(0) \in \mathfrak a$. Hence, there exists	$g\in S$ such that $ f(1) -f(0) = g(1) - g(0).$ Rearranging this equality, we get
		\[ (f-g)(0)=f(0)-g(0)=f(1)-g(1)=(f-g)(1)\in Z_{(2)}  \text{ since } f,g\in R_o\]
		Therefore $f-g\in R\subseteq S$. Since $g\in S$, $f=g+(f-g)\in S$, proving the claim.
	
		If $\mathfrak a$ we equal to $\mathbb Z_{(2)}$, then every $f\in R_o$ would satisfy $f(1)-f(0)\in\mathfrak a$ since both $f(0)$ and $f(1)$ belong to $\mathbb Z_{(2)}$. But then $S=R_o$, contrary to $S\subsetneq R_o$. Thus, $\mathfrak a$ is a proper ideal of $\mathbb Z_{(2)}$. Hence, $\mathfrak a=0$ or $\mathfrak a=2^n\mathbb Z_{(2)}$ for some 		$n\geq1$. Thus, $S=R$ or
		\[ S= \{f\in R_o: f(1)-f(0)\in2^n\mathbb Z_{(2)}\}.	\]
	
		Now let $f\in S$. In either case, $f(0)$ and $f(1)$ have the same image modulo $2\mathbb Z_{(2)}$. If $f(0)\notin2\mathbb Z_{(2)}$, both values are units of $\mathbb Z_{(2)}$. Since $f(0)\neq0$ and $f(1)\neq0$, the reciprocal $f^{-1}=1/f$ belongs to $W$. Its values $1/f(0)$ and $1/f(1)$ belong to $\mathbb Z_{(2)}$, so $f^{-1}\in R_o$. Furthermore,
	
		\[	f^{-1}(1)-f^{-1}(0)	=-\frac{f(1)-f(0)}{f(0)f(1)}= \underbrace{\left(-\frac{1}{f(0)f(1)}\right)}_{\in\mathbb Z_{(2)}}
		\underbrace{\bigl(f(1)-f(0)\bigr)}_{\in\mathfrak a}	\in\mathfrak a.	\]
		Therefore, $f^{-1}\in S$. Suppose that $f(0)\in2\mathbb Z_{(2)}$. If $f$ were a unit of $S$, there would exist $g\in S$ such that $fg=1$. Since $g\in S\subseteq R_o$, we have $g(0)\in\mathbb Z_{(2)}$, and therefore $1=(fg)(0)=f(0)g(0)\in2\mathbb Z_{(2)},$ a contradiction. Thus, $f$ is not a unit of $S$. Consequently, the nonunits of $S$ form the ideal $\{f\in S:f(0)\in2\mathbb Z_{(2)}\}.$ Thus, every such $S$, including $R$, is local. Finally, consider the overring $H=\{f\in W:f(0)=f(1)\}$ of $R$. We have $X\in R_o\setminus H,$ and $1/2\in H\setminus R_o.$ Hence, $R_o$ and $H$ are incomparable. Thus $\overline R$ is Prüfer, every intermediate ring between $R$ and $\overline R$ is comparable to $R_o=\overline R$, and every intermediate ring properly below $R_o$ is local.	Nevertheless, $R_o$ is not comparable. Therefore ``normally closed'' cannot be replaced by ``quasilocal'' in \cite[Theorem 24]{ayache}.
	\end{example}
	We now consider \cite[Problem 30]{ayache}. This question asks whether $\overline R$ must be Prüfer when a quasilocal domain $R$ has a comparable overring $R_o$ of Type III. In this case, $R_o$ is integral over $R$ and has a unique maximal ideal $Q_o$ properly containing the maximal ideal $\mathfrak m$ of $R$. The following example shows that $\overline R$ need not be Prüfer, even when $R_o$ is the minimal overring of $R$. For the next example, let $X,Y$ be indeterminates over $\mathbb Q$, and put
	\[ W= \left\{ \frac{p(X,Y)}{q(X,Y)}: p,q\in\mathbb Q[X,Y],\ q(X,0)\neq0\text{ in }\mathbb Q[X] \right\}. \]
	Thus, $W=\mathbb Q[X,Y]_{(Y)}$ is a local domain with maximal ideal $YW$ and quotient field $\mathbb Q(X,Y)$. Since every nonzero ideal of $W$ has the form $Y^nW$ for some $n\geq0$, $W$ is a principal ideal domain, and hence a discrete valuation domain.

	Every $f\in W$ has a unique expression $f=a(X)+Yb(X)+Y^2h(X,Y) $, where $a,b\in\mathbb Q(X)$ and $ h\in W. $ For instance,
	\[ \frac{1}{1-Y}=1+Y+Y^2\frac{1}{1-Y}, \]
	where $a(X)=b(X)=1$ and $h(X,Y)=1/(1-Y)\in W$.

	For $b\in\mathbb Q(X)$, we define $\lambda(b)\in\mathbb Q$ to be the constant coefficient of the formal Laurent expansion of $b$ in powers of $X$. For example,
	\[ \frac{1}{X(1-X)} =X^{-1}+1+X+X^2+\cdots, \qquad \lambda\left(\frac{1}{X(1-X)}\right)=1. \]
	The map $\lambda:\mathbb{Q}(X)\longrightarrow \mathbb{Q}$ is $\mathbb Q$-linear, and $\lambda(1)=1,\hspace{2mm}$ $ \lambda(X^j)=0$ for every $j\in\mathbb Z\setminus\{0\}.$

\begin{example}\label{counterexample-problem-30}
	Using the notation above, define
	\[ R_o= \{a+Yb+Y^2h: a\in\mathbb Q,\ b\in\mathbb Q(X),\ h\in W\} \hspace{2mm}\text{ and }\hspace{2mm}R=\{a+Yb+Y^2h\in R_o: \lambda (b) =0\}.\]
	We have $R_o=\mathbb Q+YW$. Clearly, $R_o\subseteq\mathbb Q+YW$. Conversely, let $a+Yg\in\mathbb Q+YW$, where $a\in\mathbb Q$ and $g\in W$. Set $b(X):=g(X,0)\in\mathbb Q(X)$.
	Since $g-b(X)\in YW$, write $g=b(X)+Yh$ for some $h\in W$. Then $a+Yg=a+Yb(X)+Y^2h\in R_o.$ Moreover, 
	\[ YX^j\in R \quad(j\in\mathbb Z\setminus\{0\}), \qquad Y^2W\subseteq R, \qquad Y\in R_o\setminus R. \]
	
	One can verify that $R$ is a subring of $R_o$ by using linearity of $\lambda$. Both rings have the quotient field $\mathbb Q(X,Y)$.
	Indeed, for every $f\in W$,
	\[ f=\frac{Y^2f}{Y^2}\in T(R) \]
	since $Y^2f,Y^2\in R$ and $Y^2\neq0$.
	Since $R\subseteq R_o\subseteq W$,
	\[  T(R) = T(R_o) = T(W) =\mathbb Q(X,Y). \]
	
	\smallskip
	We next show that $R$ and $R_o$ are local.
	Consider the homomorphism $\varphi:R_o\longrightarrow\mathbb Q $ defined by $f\longmapsto f(X,0).$ Since $R\subseteq R_o$, its restriction gives a homomorphism $\varphi|_R: R \longrightarrow \mathbb Q$. Every rational number belongs to $R$ as a constant function, so both maps are surjective. Their kernels are $\ker\varphi=YW:=Q_o$ and $\ker(\varphi|_R)=\{f\in R:f(X,0)=0\}:=\mathfrak m.$ The First Isomorphism Theorem gives $R_o/Q_o\cong \mathbb Q$ and $R/\mathfrak m\cong\mathbb Q$. Hence, $Q_o$ and $\mathfrak m$ are maximal ideals of $R_o$ and $R$, respectively.
	
	Let $f=a+Yb(X)+Y^2h(X,Y)\in R_o$, where $a\in\mathbb Q\setminus\{0\}$. Since $f(X,0)=a\neq0$, we have $f\notin YW$, so $f$ is a unit of $W$. Therefore, $f^{-1}$ is of the form $f^{-1}=c(X)+Yd(X)+Y^2k(X,Y)$, where $c(X),d(X)\in\mathbb Q(X)$ and $k(X,Y)\in W$. Multiplying both sides by $f$ gives
	\[1=ff^{-1}=ac(X)+Y\bigl(ad(X)+b(X)c(X)\bigr)+Y^2L(X,Y)\]
	for some $L(X,Y)\in W$. We set $Y=0$ so that $1=ac(X)$, and hence $c(X)=a^{-1}.$ On the other hand, we have $0=ad(X)+b(X)c(X)+YL(X,Y)$ implying that $ad(X)+b(X)c(X)=0.$ We substitute $c(X)=a^{-1}$ and solve for $d(X)$ so that $d(X)=-a^{-2}b(X).$ Consequently,
	\[ f^{-1}=a^{-1}-a^{-2}Yb(X)+Y^2k(X,Y). \]
	Since $a^{-1}\in\mathbb Q$, $f^{-1}\in R_o$. On the other hand, if $f\in R$, then $\lambda(b(X)) =0$. As $-a^{-2}\in\mathbb Q$ and $\lambda$ is $\mathbb Q$-linear, we have
	\[ \lambda\bigl(-a^{-2}b(X)\bigr) =-a^{-2}\lambda(b(X))=0. \]
	Thus, $f^{-1}\in R$ as well, so that every element of $R_o\setminus Q_o$ is a unit of $R_o$, and every element of $R\setminus\mathfrak m$ is a unit of $R$. Consequently, $Q_o$ and $\mathfrak m$ are the unique maximal ideals of $R_o$ and $R$, respectively. Moreover, we have $Q_o\cap R=\mathfrak m$ and $ Y\in Q_o\setminus\mathfrak m$, and hence $\mathfrak m\subsetneq Q_o$.\looseness=-1
	
	\smallskip
	We now show that $R_o=\overline R$.
	Let $f(X,Y)=a+Yb(X)+Y^2h(X,Y)\in R_o$, and put
	$c=\lambda(b(X))\in\mathbb Q$.
	Since $\lambda(b(X)-c)=0$, we have
	\[
	g(X,Y):=f(X,Y)-cY=a+Y(b(X)-c)+Y^2h(X,Y)\in R.
	\]
	This implies that $f=g+cY\in R+\mathbb{Q}Y$. Thus, $R+\mathbb QY=R_o=R[Y]$. As $Y^2\in R$, $Y $ is the root of the monic polynomial $Z^2-Y^2\in R[Z]$. Therefore, $R_o$ is integral over $R$.
	
	\smallskip To show that $R_o$ is integrally closed, let $z\in\mathbb Q(X,Y)$ be integral over $R_o$. Then, there exists a monic polynomial 
	\[F(Z)=Z^n+c_{n-1}(X,Y)Z^{n-1}+\cdots+c_0(X,Y)\in R_o[Z]\] 
	such that $F(z)=0.$	Since $R_o\subseteq W$, all coefficients of $F$	belong to $W$. Thus ,$F\in W[Z]$ is a monic polynomial with root $z$, so $z$ is integral over $W$. The valuation domain $W$ is integrally closed, and hence $z\in W$. By the definition of $R_o$, we have $a_i:= c_i(X,0)\in\mathbb Q$ for $0\leq i\leq n-1.$ Setting $Y=0$ in $F(z)=0$ we get $z(X,0)^n+ a_{n-1}z(X,0)^{n-1}+\cdots+a_0=0.$ Therefore, $z(X,0)\in\mathbb Q(X)$ is a root of the monic polynomial
	\[ P(Z)=Z^n+a_{n-1}Z^{n-1}+\cdots+a_0 \in\mathbb Q[Z]. \]
	Thus, $z(X,0)$ is algebraic over $\mathbb Q$. The only elements of $\mathbb Q(X)$ algebraic over $\mathbb Q$ are the elements of $\mathbb Q$. Consequently, $z(X,0)\in\mathbb Q$. Since $z\in W$, we obtain $z-z(X,0)\in YW,$ and hence $z\in\mathbb Q+YW=R_o$. Thus, $R_o$ is integrally closed. Since $R_o$ is also integral over $R$, we conclude that $\overline R=R_o$.
	
	We next prove that every proper overring of $R$ contains $R_o$. First observe that if $s\in R_o\setminus R$ then  $s=r+cY$ for some $r\in R$ and $0\neq c\in\mathbb Q$. Since $c^{-1}\in\mathbb Q\subseteq R$, we obtain $Y=c^{-1}(s-r)\in R[s].$ Therefore, $R_o=R[Y] \subseteq R[s]$.
	
	Let $T$ be a proper overring of $R$, and choose $t\in T\setminus R$. If $t\in R_o$, the preceding observation gives $R_o\subseteq R[t]\subseteq T$.
	
	Suppose next that $t\in W\setminus R_o$. Since $t\notin R_o$, $t(X,0)$ does not belong to $\mathbb Q$. Put $c(X):=t(X,0)$. Thus, $c(X)$ is a nonconstant element of $\mathbb Q(X)$. Write its Laurent expansion at $X=0$ as $$c(X)=\sum_{k\geq k_0}c_kX^k.$$ Since $c(X)$ is nonconstant, there exists $j\neq0$ such that $c_j\neq0$. Put $r:=YX^{-j}\in R.$ Since $t(X,0)=c(X)$, we have $t-c(X)\in YW$. Thus, $t=c(X)+Yg(X,Y)$ for some $g(X,Y)\in W$. Using $r=YX^{-j}$, we obtain \[tr=\bigl(c(X)+Yg(X,Y)\bigr)YX^{-j}=Y\bigl(c(X)X^{-j}\bigr)+Y^2\bigl(g(X,Y)X^{-j}\bigr).\]
	Since $X^{-j}\in W$, we have $g(X,Y)X^{-j}\in W$. Thus $(tr)(X,0)=0$, and the coefficient of $Y$ in this decomposition is $c(X)X^{-j}$. Since $\lambda\bigl(c(X)X^{-j}\bigr)=c_j\neq0,$ we have $tr\in R_o\setminus R$. Applying the preceding observation to $tr\in T$, we get $R_o\subseteq T$.
	
	Finally, suppose that $t\notin W$. Since $W$ is a discrete valuation domain with maximal ideal $YW$, we write $t=Y^{-d}u$, where $ d\geq1 $ and $ u\in W^\times.$ Put $\ell(X) :=u(X,0)\neq0$, and set $\rho:=Y^{d+1}/\ell(X).$ As $1/\ell(X)\in W$ and $d+1\geq2$, we have $\rho\in Y^2W\subseteq R$. Moreover, since $\ell(X)=u(X,0)$, we can write $u(X,Y)=\ell(X)+Yg(X,Y)$ for some $g(X,Y)\in W$. Hence
	\[	\rho t	=Y\frac{u(X,Y)}{\ell(X)} =Y\frac{\ell(X)+Yg(X,Y)}{\ell(X)}	=Y+Y^2\frac{g(X,Y)}{\ell(X)}.	\]
	Since $g(X,Y)/\ell(X)\in W$, $\rho t\in R_o\setminus R$, and again $R_o\subseteq T$. Therefore, every proper overring of $R$ contains $R_o$. In particular, $R_o$ is a minimal comparable overring of $R$. Since $\mathfrak m\subsetneq Q_o$, it is of the third type. Since $Y^{-1}\notin W$, we have $Y^{-1}\notin R_o$. So, $R\subsetneq R_o\subsetneq\mathbb Q(X,Y).$ However, $X\notin R_o$ and $X^{-1}\notin R_o$ because their values at $Y=0$ are not elements of $\mathbb Q$. Hence, $R_o$ is not a valuation domain. Since a local Prüfer domain is a valuation domain, $\overline{R}=R_o$ is not Prüfer.
\end{example}
\section*{Acknowledgments}

The results presented in this paper constitute part of
the second author's Ph.D. thesis.
The second author acknowledges the financial support of
T{\"U}B{\.I}TAK through the B{\.I}DEB 2211 National
Ph.D. Scholarship Program.
\begin{thebibliography}{99}
	\bibitem{marot} B. Ay Saylam and E. Canan, \emph{Comparable overrings of Marot rings}, Journal of Algebra and Its Applications \textbf{25}(9) (2026), 2750161, DOI: \href{https://doi.org/10.1142/S0219498827501611}{10.1142/S0219498827501611}.
	
	
	\bibitem{ayache} A. Ayache, \emph{When is a fixed ring comparable to all overrings?}, Journal of Algebra and Its Applications \textbf{21} (2022), 2250050, DOI: \href{https://doi.org/10.1142/S0219498822500505}{10.1142/S0219498822500505}.
	
	
	\bibitem{huckaba} J. A. Huckaba, \emph{Commutative Rings with Zero Divisors}, Marcel Dekker, New York, 1988.
	
	
	\bibitem{atiyah} M. F. Atiyah and I. G. Macdonald, \emph{Introduction to commutative algebra}, CRC Press, 2018.
	
	
	\bibitem{gilmer-heinzer} R. Gilmer and W. J. Heinzer, \emph{Intersections of quotient rings of an integral domain}, Journal of Mathematics of Kyoto University \textbf{7}(2) (1967), 133--150.
	
	
	\bibitem{normal-pairs} M. Ben Nasr and N. Jarboui, \emph{New results about normal pairs of rings with zero-divisors}, Ricerche di Matematica \textbf{63} (2014), 149--155, DOI: \href{https://doi.org/10.1007/s11587-013-0169-1}{10.1007/s11587-013-0169-1}.
	
	\bibitem{chang} 
	Chang, Gyu Whan. "Generalizations of Pseudo-Valuation Rings." Trends in Commutative Rings Research (2004): 61.
	
	\bibitem{dobbs-jarboui}
	D. E. Dobbs and N. Jarboui, \emph{Prüfer-closed extensions and FCP $\lambda$-extensions of commutative rings}, Palestine Journal of Mathematics \textbf{11}(3) (2022), 362--378.
	
\end{thebibliography}

\end{document}
